% concatenated from Boyd-Mahler-AI-arXiv.tex
\documentclass[12pt]{amsart}%
\usepackage{verbatim}
\usepackage{latexsym}
\usepackage{amsmath}
\usepackage{amsfonts}
\usepackage{float}
\usepackage{tikz}
\usepackage{placeins}%
\usetikzlibrary{positioning}
\newcommand{\countpct}[2]{%
  #1\,{\scriptsize(#2\%)}%
}
\usepackage{booktabs}
\usepackage{bbm}

\usepackage{tabularx,multirow,booktabs,makecell,amssymb}
\usepackage{siunitx}
\usepackage{amssymb,hyperref}
\usepackage[normalem]{ulem}
\pagestyle{myheadings}
\textheight=23cm
\textwidth=17.7cm
\topmargin=-0.7cm
\oddsidemargin=-.64cm
\evensidemargin=-.64cm
\renewcommand{\theequation}{\arabic{equation}}

\DeclareMathOperator{\arcsinh}{arcsinh}
\usepackage{graphicx}
\usepackage{caption}
\usepackage{subcaption}

\renewcommand{\pmod}[1]{\hspace{0.05em}(\mathrm{mod}\,#1)}


\usepackage{dsfont}
\newcommand{\1}{\mathds{1}}

\renewcommand{\pmod}[1]{\,{\rm mod}\,#1}
\newcommand{\D}{\mathbb{D}}

\newcommand{\C}{\mathbb{C}}
\newcommand{\F}{\mathbb{F}}
\newcommand{\R}{\mathbb{R}}
\newcommand{\Q}{\mathbb{Q}}
\newcommand{\Z}{\mathbb{Z}}
\newcommand{\PP}{\mathbb{P}}
\newcommand{\TT}{\mathbb{T}}%
\newcommand{\HH}{\mathbb{H}}
\newcommand{\tr}{\mathrm{tr}}%
\newcommand{\E}{\mathbb{E}}
\newcommand{\dd}{\;\mathrm{d}}%


\newtheorem{thm}{Theorem}[section]
\newtheorem{defn}[thm]{Definition}
\newtheorem{prop}[thm]{Proposition}
\newtheorem{coro}[thm]{Corollary}
\newtheorem{lem}[thm]{Lemma}
\newtheorem{conj}[thm]{Conjecture}
\newtheorem{ex}[thm]{Example}
\newtheorem{qu}[thm]{Question}
\newtheorem{pa}[thm]{Possible answer}
\newtheorem{exs}[thm]{Examples}
\newtheorem{obs}[thm]{Observation}

\theoremstyle{definition}
\newtheorem{rem}[thm]{Remark}

\numberwithin{equation}{section}
\numberwithin{table}{section}
\numberwithin{figure}{section}

\newcommand{\e}{\mathrm{e}}%
\newcommand{\Res}{\mathrm{Res}}%
\newcommand{\m}{\mathrm{m}}

\newcommand{\Li}{\mathrm{Li}}%
\newcommand{\I}{\mathrm{I}}%
\newcommand{\re}{\mathop{\mathrm{Re}}}%
\newcommand{\im}{\mathop{\mathrm{Im}}}%
\newcommand{\hypgeo}[2]{%
  {\vphantom{F}}_{#1}\kern-\scriptspace F_{#2}%
}


\newmuskip\pFqmuskip
\newcommand*\pFq[6][8]{%
  \begingroup%
  \pFqmuskip=#1mu\relax
  \mathchardef\normalcomma=\mathcode`,
  \mathcode`\,=\string"8000
  \begingroup\lccode`\~=`\,
  \lowercase{\endgroup\let~}\pFqcomma
  {}_{#2}F_{#3}{\left(\!\left.\genfrac..{0pt}{}{#4}{#5}\right|#6\!\right)}%
  \endgroup
}
\newcommand{\pFqcomma}{{\normalcomma}\mskip\pFqmuskip}


\begin{document}
\raggedbottom%

\title{Machine learning the arithmetic of Boyd's Mahler measure conjectures}
\author{Alberto Alfarano}
\author{Pablo Bianucci}

\author{Matilde N. Lal\'in}
\author{Berend Ringeling}


\address{Alberto Alfarano: Axiom Math}\email{alberto@axiommath.ai}

\address{Pablo Bianucci:  Department of Physics, Concordia University, 7141 Sherbrooke St. W.
Montreal, QC  H4B 1R6, Canada} \email{pablo.bianucci@concordia.ca}


\address{Matilde Lal\'in:  D\'epartement de math\'ematiques et de statistique, Universit\'e de Montr\'eal. CP 6128, succ. Centre-ville. Montreal, QC H3C 3J7, Canada}\email{matilde.lalin@umontreal.ca}
\address{Berend Ringeling:  D\'epartement de math\'ematiques et de statistique, Universit\'e de Montr\'eal. CP 6128, succ. Centre-ville. Montreal, QC H3C 3J7, Canada}\email{bjringeling@gmail.com}

\begin{abstract}
Boyd conjectured that the Mahler measure of $P_k(x,y)=x+y+\frac{1}{x}+\frac{1}{y}+k$ for $k$ an integer, is given by $r_kL'(E_k,0)$, where $E_k$ is the elliptic curve associated to the zero locus of $P_k$ and $r_k$ is a rational number. We study various arithmetic properties of $r_k$ using a dataset containing the first $250{,}000$ values of $k$, combining
large-scale statistical analysis assisted by \textsc{Claude} with transformer-based
experiments carried out using \textsc{Axolver}.

We recover Boyd's observation that, apart from a few exceptions, $r_k$ is
the reciprocal of an integer. The size of this integer is governed by the conductor of the elliptic curve. Moreover, its $p$-adic valuations display markedly different behavior according
to the prime. For $p\geq 5$, the probability of
$v_p(r_k)=-m$
for $m\geq 1$ appears to be $p^{-m}$. For the primes  $2$ and $3$, however, we find additional arithmetic structure
involving congruence conditions on $k$ and the primes of bad reduction of
$E_k$. Although the neural networks do not predict $r_k$ exactly, they recover significant information about its magnitude and valuations. In particular, the experiments at the prime $2$ suggest arithmetic structure beyond the explicit predictor obtained from our statistical analysis.


\end{abstract}

\maketitle

The (logarithmic) Mahler measure of a nonzero rational function $P\in \C(x_1,\dots,x_n)$ is defined by
\begin{align*}
\mathrm{m}(P)=\frac{1}{(2\pi i)^n}\int_{\mathbb{T}^n} \log\left| P(x_1,\dots,x_n)\right|\frac{\dd x_1}{x_1}\cdots \frac{\dd x_n}{x_n},
\end{align*}
where $\mathbb{T}^n=\{(u_1,\dots,u_n)\in \C^n\, : \, |u_1|=\cdots=|u_n|=1\}$.
In the one variable polynomial case given by  $n=1$, Jensen's formula leads to an expression in terms of the roots of $P$. More precisely, if $P(x)= a \prod_j(x-\alpha_j)$, one can see that $\m(P)=\log |a| + \sum_j \log\max\{1,|\alpha_j|\}$.


However, for $n>1$, the nature of the values is more mysterious. In many cases where $P$ has rational coefficients, the Mahler measure has been related to values of functions with arithmetic significance.

Deninger~\cite{Deninger} predicted a connection of certain two-variable Mahler measures with special values of $L$ functions associated to elliptic curves via Be\u\i linson's conjectures, which are very general frameworks relating special values of $L$-functions to regulators, generalizing central statements such as the Birch and Swinnerton Dyer conjecture and the Dirichlet class number formula. These connections were verified by Boyd~\cite{Bo98} and Rodriguez-Villegas~\cite{RV} in many cases, most of them numerical, with a few formulas proven by Brunault, Mellit, Rogers and others \cite{Br06,La10,LR07,Me12,RZ12,RZ14,Zu14}.
One of the most studied families fitting the above framework is the one proposed by Boyd~\cite{Bo98}, given by
\[\mu_k:=\mathrm{m}\left(x+y+\frac{1}{x}+\frac{1}{y}+k\right),\]
where $k$ is an integral parameter.  Boyd computed $\mu_k$ for $k\in\Z\cap [0,40]$ and combined these observations with Deninger's predictions to conjecture that for  $k\in \Z\setminus\{-4,0,4\}$,
\begin{equation}\label{eq:conjecture}
\mu_k=r_k L'(E_k,0),
\end{equation}
where $E_k$ is the elliptic curve corresponding to the projective closure of the desingularization of
$x+y+\frac{1}{x}+\frac{1}{y}+k=0$, and $r_k$ is a nonzero rational number. In fact, apart from a few exceptions, it seems that $1/r_k$ is an integer. Table \ref{table-conj} shows $1/r_k$ for the first 40 values of $k$. The conductor $N_k$ is included in the table, since it determines the $L$-function.%
Moreover, the sign of $r_k$ should be determined by the root number $\omega_k$ of $E_k$, as the Mahler measure is positive. We remark, by using the change of variables $x\rightarrow -x$ and $y \rightarrow -y$ in the integral, that $\mu_k=\mu_{-k}$, and therefore it suffices to consider positive $k$.%
When $k=0, 4$, the zero loci of the corresponding Laurent polynomials have genus 0. In these cases, $\mu_0=0$ and $\mu_4=2L'(\chi_{-4},-1)$, the Dirichlet series on the character of conductor 4 \cite{Bo98}.


\begin{table}[h]
\centering
\begin{tabular}{ccc||ccc}
\hline
$k$ & $1/r_k$ & $N_k$ & $k$ & $1/r_k$ & $N_k$ \\
\hline
1  & 1     & 15   & 21 & $-12$  & 1785  \\
2  & 1     & 24   & 22 & 24     & 3432  \\
3  & 1/2   & 21   & 23 & 6      & 1311  \\
4  & --    & --   & 24 & 8      & 840   \\
5  & 1/6   & 15   & 25 & $-16$  & 3045  \\
6  & 2     & 120  & 26 & 128    & 17160 \\
7  & 2     & 231  & 27 & 12     & 2139  \\
8  & 1/4   & 24   & 28 & $-2$   & 336   \\
9  & 2     & 195  & 29 & $-24$  & 4785  \\
10 & $-8$  & 840  & 30 & 240    & 26520 \\
11 & $-8$  & 1155 & 31 & $-16$  & 3255  \\
12 & 1/2   & 48   & 32 & 1/3    & 42    \\
13 & $-4$  & 663  & 33 & $-204$ & 35409 \\
14 & 8     & 840  & 34 & 256    & 38760 \\
15 & $-24$ & 3135 & 35 & 224    & 42315 \\
16 & 1/11  & 15   & 36 & 2      & 240   \\
17 & $-24$ & 4641 & 37 & $-208$ & 50061 \\
18 & $-16$ & 1848 & 38 & 288    & 54264 \\
19 & $-40$ & 6555 & 39 & 336    & 58695 \\
20 & 2     & 240  & 40 & $-8$   & 1320  \\
\hline
\end{tabular}
\caption{Numerical values of $r_k$ and the conductor of $E_k$.\label{table-conj}}
\end{table}


Some of these identities have been proven, see Table \ref{table:proven_integer}.
\begin{table}[h]
\centering
\begin{tabular}{cccc}
\hline
$k$ & $1/r_k$ & Conductor of $E_k$ & Reference \\
\hline
1  & 1      & 15   & Rogers and Zudilin \cite{RZ14} \\
2  & 1      & 24   & Rogers and Zudilin \cite{RZ12} \\
3  & 1/2    & 21   & Brunault \cite{Brunault-Siegel}; Lal\'in, Samart, and Zudilin \cite{LSZ16} \\
5  & 1/6    & 15   &  Lal\'in \cite{La10} from $k=1$ \\
8  & 1/4    & 24   & Lal\'in and Rogers \cite{LR07} from $k=2$\\
12 & 1/2    & 48   & Brunault \cite{Brunault-Siegel} \\
16 & 1/11   & 15   & Lal\'in \cite{La10} from $k=1$ \\
\hline
\end{tabular}
\caption{Proven formulas for integer values of $k$.}
\label{table:proven_integer}
\end{table}


Understanding the nature of these identities would shed light on our understanding of special values of $L$ functions and Be\u\i linson's conjectures. The following is a natural question to pose.
\begin{qu}
    How does $r_k$ relate to $k$ and to the elliptic curve $E_k$?
\end{qu}

We remark that $r_k$ should be predicted by the Bloch--Kato conjectures. Very recent work by Dummigan, Golyshev, de Jeu and Kerr \cite{DGJK} investigates the $2$-part of the Bloch--Kato conjecture for this family of curves.


We have computed the first $250{,}000$ values of $\mu_k$ and explored this question in two different ways: by analyzing the data with the assistance of  \textsc{Claude} and by training neural networks on the  problem via \textsc{Axolver}%
, with the goal of predicting certain properties of $r_k$.

The data analysis
revealed several arithmetic regularities
exhibited by $r_k$. After accounting for the sign, which is governed by the root number, the reciprocal $n_k=1/r_k$ is almost always an integer. Its order of magnitude is explained by the conductor through the estimate
\[
        |n_k|\asymp \frac{N_k}{\log (k)},
\]
while its $p$-adic valuations show different behaviors depending on the prime. For primes $p\geq 5$, the data suggest the simple law
\[
        \mathbb{P}(v_p(n_k)=r)\stackrel{?}{=}p^{-r}\qquad (r\geq 1).
\]
The primes $2$ and $3$, however, have a different behavior. The $3$-adic valuation seems to be strongly influenced by congruence conditions on $k$ and by the parity of the number of bad primes congruent to $1\pmod 3$. The $2$-adic valuation appears to be  well modeled by a deterministic term depending on primes of bad reduction plus an apparently random negative-binomial error term.%
More precisely, we observe that
\[v_2(n_k) \overset{?}{=} \omega_{\mathrm{odd}}(k) + 2\,\omega_{\mathrm{odd}}(k^2-16) - c(k) + s(k) + X_k,\]
where $\omega_{\mathrm{odd}}(n)$ counts the odd primes dividing $n$, and
$c(k)$ and $s(k)$ are small bounded correction terms that codify the behavior of $E_k$ at $2$, defined in \eqref{eq:c-defn} and \eqref{eq:s-defn}, respectively.
Finally, the empirical distribution of the residual $X_k$ is well approximated by the negative-binomial distribution
 $\mathrm{NB}\!\left(r=3,\;p=\tfrac{3}{5}\right)$. We emphasize that this is a heuristic observation rather than a precise conjecture. The work of Dummigan, Golyshev, de Jeu and Kerr \cite{DGJK}
provides a formula for $v_2(n_k)$ for an infinite, positive-density
family of $k$, conditionally on the Bloch--Kato conjecture. The formula
depends on the 2-valuation of a Selmer group associated to an elliptic curve 2-isogenous to $E_k$
as well as on a regulator index
$\iota_{k/4,2}\in\{1,2\}$, which measures the index of an explicit
motivic cohomology class in the rank-one $\Z_{(2)}$-lattice determined
by the $2$-adic regulator. We briefly examine this connection in Section \ref{ssec:lowboundproof}.


The experiments showed us that \textsc{Axolver} is able to recover part of this arithmetic structure from the data, but that the full problem of predicting $r_k$ remains very difficult. While the model is unable to completely learn $r_k$, it  does learn meaningful information about its size, with the best behavior occurring when the conductor is included among the input features. In particular, the model begins to reproduce the observed distribution of $L(E_k,2)$. The valuation experiments are more mixed. For $p=5$ and $p=7$, the model tends to collapse to the most frequent output $v_p(n_k) = 0$, so these experiments do not yet reveal useful structure. By contrast, the experiments for $v_2(n_k)$ and $v_3(n_k)$ show that the network can exploit arithmetic features such as the factorization of the discriminant or conductor, giving evidence that the statistical patterns observed in Section~\ref{sec:statistics} are visible to a learning model. In the case of the prime $2$, the model's learning seems to go beyond the understanding that we can derive from the data  (see Section \ref{sec:secondplateau}).

This article is organized as follows.
Section \ref{sec:background} presents the general mathematical background for the problem. Section \ref{sec:data-generation} briefly discusses the generation of the data. The statistical observations are presented in Section \ref{sec:statistics}. Section \ref{sec:model} describes the model and the experiments run in \textsc{Axolver}, while  Section \ref{sec:machine} contains a detailed discussion on the results of those experiments. Section \ref{sec:conclusions} summarizes our findings and discusses directions for further research, while Section \ref{sec:appendix} contains an appendix with supplementary tables and figures.

\section*{Acknowledgments} This project was initiated during the workshop ``Using AI to find examples'' that took place in April 2026 at the Centre de recherches math\'ematiques. The authors are grateful to the organizers, particularly Andrew Granville, and to the CRM for the excellent working conditions. The authors are thankful to Fran\c cois Charton, Christopher Deninger, Jordan Ellenberg, Dimitris Koukolopoulos,  Shousen Lu, Riccardo Pengo, and Wadim Zudilin for several helpful discussions. Part of this work was conducted during the workshop ``AI and number theory" that took place in May 2026 at the American Institute of Mathematics. The authors thank the organizers and Anthropic for providing access and guidance to \textsc{Claude Max} during this event.

This work was partially supported by the Natural Sciences and Engineering Research Council of Canada (RGPIN-2019-06988 to PB, RGPIN-2022-03651 to ML), and the Fonds de recherche du Qu\'ebec - Nature et technologies (Projet de recherche en \'equipe 345672 to ML), the  Institut des sciences math\'ematiques (postdoctoral fellowship to BR),  the Centre de recherches math\'ematiques (postdoctoral fellowship to BR).


\section*{Declaration of Generative AI and AI-Assisted Technologies} During the preparation of this manuscript, the authors used artificial intelligence (AI) tools to assist with data analysis, machine learning modeling, and manuscript editing:

\begin{itemize}\item Data Analysis: Anthropic's \textsc{Claude} model was utilized to process, clean, and identify patterns within the research dataset.

\item Machine Learning: Axiom's \textsc{Axolver} framework was used to train sequence-to-sequence neural networks and execute mathematical/symbolic reasoning tasks within the study's machine learning pipeline.

\item Writing and Editing: OpenAI's \textsc{ChatGPT} was used  to improve the readability, grammar, and style of the author-drafted text.
\end{itemize}

The original research framework, data collection, and final interpretations were entirely conceptualized and directed by the human authors. After using these tools, the authors reviewed, verified, and edited all outputs to ensure accuracy and eliminate potential hallucinations. The authors accept full accountability for the content and integrity of the final work.


\section*{Data Availability Statement}
The code used in this study is available in the \textsc{GitHub} repository
\url{https://github.com/BerendRingeling/Mahler}. The datasets are available
through the \textsc{Dropbox} link provided in that repository.


\section{Background and generalities} \label{sec:background}

The plane curve defined by the equation
\[x+\frac{1}{x} + y + \frac{1}{y} + k=0,\]
where $k$ is a parameter, is birationally equivalent to an elliptic curve (over $\Q(k)$). This can be seen by applying, for example, the birational change of variables
\[
x=\frac{8(k^2X+8Y)}{k(k^2-16X)},
\qquad
y=\frac{8(k^2X-8Y)}{k(k^2-16X)},
\]
with inverse
\[
X=-\frac{k^2xy}{16},
\qquad
Y=\frac{k^3(xy+1)(x-y)}{128},
\]
which
leads to the Weierstrass equation
\begin{equation}\label{eq:Weierstrass}
E_k: Y^2 = X\left(X^2 + \frac{k^2}{8}\left(\frac{k^2}{8} - 1\right)X + \frac{k^4}{16^2}\right).
\end{equation}
For fixed $k$, this curve is nonsingular provided that $k\neq 0,\pm4$.


If an elliptic curve $E$ is defined over $\Q$, one can construct its $L$-function as follows
\[L(E,s) = \prod_{\mathrm{good}\, p} (1-a_p p^{-s}+p^{1-2s})^{-1} \prod_{\mathrm{bad}\, p} (1-a_p p^{-s})^{-1}=\sum_{n=1}^\infty \frac{a_n}{n^s},\]
where for $p$ prime,
\[a_p = 1+p-\#E(\F_p)\]
denotes the trace of Frobenius.
Here the bad $p$ are the primes where the corresponding reduction of $E$ modulo $p$ is singular, and the good $p$ are all the other primes. The bad primes are exactly those that divide the conductor $N$ of the elliptic curve.

The Weierstrass model \eqref{eq:Weierstrass} is defined over $\Q(k^2)$. This observation prompted Rodriguez-Villegas \cite{RV} to conjecture that equation \eqref{eq:conjecture} extends to the case where $k^2$ is a positive integer and to prove the cases $k=2\sqrt{2}, 4\sqrt{2}$. A word of caution here is that, since $k \not \in \Q$, and any change of variables will be defined over $\Q(k)$, the choice of the model does make a difference. A few more cases where $k^2\in \Z$ have been proven, see Table \ref{table:proven_nonrational}.
\begin{table}[h]
\centering
\begin{tabular}{cccc}
\hline
$k$ & $1/r_k$ & Conductor of $E_k$ & Reference \\
\hline
$\sqrt{2}$     & 4  & 56  & Zudilin \cite{Zu14} \\
$2\sqrt{2}$   & 1    & 32  & Rodriguez-Villegas \cite{RV} \\
$3\sqrt{2}$   & 2/5    & 24  & Lal\'in \cite{La10} from $k=2$  \\
$4\sqrt{2}$    & 1    & 64  & Rodriguez-Villegas \cite{RV} \\
$i$            & 1/2  & 17  & Zudilin \cite{Zu14} \\
$2i$           & 1    & 40  & Mellit \cite{Mellit-Oberwolfach} \\
$3i$           & 1/5    & 15  &  Lal\'in \cite{La10} from $k=1$ \\
$i\sqrt{2}$   & 2/3    & 24  & Lal\'in \cite{La10} from $k=2$  \\
\hline
\end{tabular}
\caption{Proven formulas for quadratic irrational and non-real values of $k$.}
\label{table:proven_nonrational}
\end{table}

In the present work, however, we restrict our statistical and  machine-learning analysis to integral values of $k$. This is the setting of Boyd's original numerical experiments.
This restriction is also convenient computationally: generating data for $k=1,\ldots,2B$ appears slightly faster than carrying out the corresponding computations for all integral values $|m| \leq B$, where $m=k^2$ is no longer required to be a square. We briefly examine the nonsquare case in Section~\ref{sec:sqrtm}, where the near-integrality, and valuation laws are found to persist in many cases.


The motivation for conjecturing formulas such as \eqref{eq:conjecture} comes from Be\u\i linson's conjectures, as explained in \cite{Bo98,Deninger,RV}.


\section{Numerical generation of data}
\label{sec:data-generation}

In order to generate the data needed to train the model, we need to compute both sides of equation \eqref{eq:conjecture}.
The computation of $L'(E_k,0)$ can be performed to very high precision in
\textsc{PARI/GP} \cite{PARI}. For the Mahler measure we have the following explicit formulas for real $k$
(see \cite[p.~2324]{RZ14}):
\begin{equation}
\label{HG}
\mu_k =
\begin{cases}
\log |k| - \dfrac{2}{k^2}\,
{}_4F_3\!\left(\frac{3}{2}, \frac{3}{2}, 1, 1;\, 2, 2, 2;\, \frac{16}{k^2}\right),
& \text{if } |k| \ge 4, \\[6pt]
\dfrac{|k|}{4}\,
{}_3F_2\!\left(\frac{1}{2}, \frac{1}{2}, \frac{1}{2};\, \frac{3}{2}, 1;\, \frac{k^2}{16}\right),
& \text{if } |k| < 4 .
\end{cases}
\end{equation}
Here ${}_4F_3$ and ${}_3F_2$ denote hypergeometric functions. This explicit
description allows one to compute $\mu_k$ to high precision for
$k \in \Z \setminus\{-4, 0, 4\}$. Both $\mu_k$ and $L'(E_k,0)$ were computed using $19$ decimal digits of precision.

Using the LLL algorithm (implemented as \texttt{lindep} in \textsc{PARI/GP}
\cite{PARI}), we can then quickly numerically recognize the rational quotient $r_k$.


We computed the values $r_k$ for $1 \le k \le 250{,}000$ (with $k \neq 4$).
In total this required about \SI{9.6e7}{\second} ($\approx 3$ years) of
single-core-equivalent computation time, with the average computation time per value rising
from roughly \SI{35}{\second} for small $k$ to about \SI{16.5}{\minute} for
$k$ near 250{,}000. The computations were carried out on 21 CPU cores
across 17 servers with a combined \SI{2108}{\giga\byte} ($\approx
\SI{2.1}{\tera\byte}$) of memory.

\section{Statistical observations of the data}
\label{sec:statistics}

We analyzed the values $r_k$ for $1 \le k \le 250{,}000$. First of all, since the Mahler measure is always positive, the sign of $r_k$ is determined by the sign of $L'(E_k,0)$, which, by the functional equation of the $L$-function, is determined by the root number $\omega_k$ of $E_k$. Secondly, for all  but a few exceptional values of $k$ in the dataset,%
namely for
\begin{equation}\label{eq:intcond}
k \notin \{3,4,5,8,12,16,32\},
\end{equation}
the quantity $r_k$ is conjecturally the reciprocal of an integer.
Inspired by this, we write $n_k =  1/r_k$, so that $n_k$ is an integer for all the $k$ satisfying \eqref{eq:intcond}. We make the following empirical observations on the size of $n_k$ and its $p$-adic valuations for primes $p$ for these $k$.


\subsection{The size of $|n_k|$}
We can use the approximation for $\mu_k$ and $L'(E_k,0)$ to find an estimate for $|n_k|$.
Using the hypergeometric description for the Mahler measure \eqref{HG}, we find
\begin{equation}
\label{mmk}
\mu_k = \log(k) + o(1)
\end{equation}
as $k \to \infty$. For $L'(E_k,0)$, we use the functional equation to arrive at
\begin{equation}
\label{lvk}
L'(E_k,0) = \omega_k \frac{N_k}{4 \pi^2 }L(E_k,2),
\end{equation}
where $N_k$ is the conductor of $E_k$ and $\omega_k$ the root number. Using the following lemma, we can estimate the size of $L(E_k,2)$.

\begin{lem}
\label{lem:lvbound}
For any elliptic curve $E$ over $\Q$, we have the bounds
    \[ (0.2117\ldots =) \, \frac{\zeta(3)^2}{\zeta(3/2)^2} \leq L(E,2) \leq \zeta(3/2)^2 \, (= 6.8245\ldots).\]
\end{lem}
\begin{proof}
We write the $L$-value in terms of its local factors
\[
L(E, 2) = \prod_{\text{primes }p} L_p(E,2),
\]
where
\[
L_p(E,2) = \begin{cases}(1-a_pp^{-2}+p\cdot p^{-4})^{-1},&p\nmid N\\ (1\pm p^{-2})^{-1}
,&p\|N\\ 1,& p^2|N,\end{cases}
\]
where $N$ is the conductor of $E$.
For good primes (i.e. $p \nmid N$), we have, by the Hasse-Weil bound, $|a_p| \leq 2 \sqrt{p}$, and therefore,
\[
(1 - p^{-3/2})^2 = 1 - 2p^{-3/2} + p^{-3} \leq 1-a_pp^{-2}+p\cdot p^{-4} \leq 1 + 2p^{-3/2} + p^{-3} = (1 + p^{-3/2})^2.
\]
Taking reciprocals,
\[
\frac{1}{(1 + p^{-3/2})^2} \leq L_p(E,2) \leq \frac{1}{(1 - p^{-3/2})^2}.
\]
For $p \| N$, the same estimates hold:
\[
(1-p^{-3/2})^2 \leq 1 - p^{-2} \leq 1 \pm p^{-2} \leq 1 + p^{-2} \leq (1 + p^{-3/2})^2,
\]
where the outer inequalities can be easily verified. Finally, for $p^2 | N$, it is immediate that these bounds hold. Taking the product over all primes $p$, we conclude that
\[
\frac{\zeta(3)^2}{\zeta(3/2)^2} = \prod_{\text{primes }p} \frac{1}{(1 + p^{-3/2})^2} \leq L(E,2) \leq \prod_{\text{primes }p} \frac{1}{(1 - p^{-3/2})^2} = \zeta(3/2)^2.
\]
\end{proof}

As a consequence of \eqref{mmk}, \eqref{lvk} and Lemma \ref{lem:lvbound}, we conclude that
\begin{equation}
\label{orderofnk}
|n_k| = \frac{|L'(E_k,0)|}{\mu_k} \asymp \frac{N_k}{\log(k)}.
\end{equation}
Thus, the order of magnitude of $n_k$ is completely determined by $k$ and the conductor $N_k$. In Figure~\ref{fig:distrL_E_2} we see the distribution of $L(E_k,2)$ in our dataset.
\begin{figure}[!htbp]
    \centering
   \includegraphics[scale = 0.5]{L_Ek_2_full.pdf}
    \caption{Values of $L(E_k, 2)$ for $k \leq 250{,}000$ in our dataset.}%
    \label{fig:distrL_E_2}
\end{figure}

Although Lemma \ref{lem:lvbound} only forces $(0.2117,6.8245)$, the values we observe are in a narrower range $(0.5222, 1.7532)$, with mean value $0.8326$. The exact distribution of $L(E_k,2)$ is not known to us; here we record only its empirical behavior and its role in governing the size of $n_k$ through \eqref{orderofnk}. We do not assert any direct relation between this distribution and the $p$-adic valuations of $n_k$ considered below.


\subsection{$p$-adic observations for primes $p \geq 5$}
\label{subsection:big_p-adic}

We observe empirically that
\[
\mathbb P\bigl(v_p(n_k)\geq 1\bigr)=\frac{1}{p-1},
\]
rather than the probability \(1/p\) predicted by a naive uniform-distribution model modulo \(p\). Conditional on \(p\mid n_k\), however, the higher valuations appear to follow the usual geometric distribution:
\[
\mathbb P\bigl(v_p(n_k)=r\mid v_p(n_k)\geq 1\bigr)
=
\left(1-\frac{1}{p}\right)p^{-(r-1)},
\qquad r\geq 1.
\]
Equivalently,
\[
\mathbb P\bigl(v_p(n_k)=r\bigr)=\frac{1}{p^r},
\qquad r\geq 1.
\]
Thus, the unexpected feature is the enhanced probability \(1/(p-1)\) of divisibility by \(p\). Once \(p\mid n_k\), the distribution of the higher valuation agrees with that of a random \(p\)-adic integer conditioned to be divisible by \(p\). Experimental results for several primes \(p\) and valuations \(r\) are recorded in Table~\ref{tab:v5_v7_v11_v13_v17_v19_v23_v29_v31_v37_statistics_cumulative}.

To quantify how well our data matches with the proposed random model, we compute $\chi^2$. Concretely, for each prime $p$, we group the values $v_p(n_k)$ into finitely many bins. For instance, for $p = 5$, we use the seven bins $v_5(n_k) = r$ for $0 \leq r \leq 5$ and $v_5(n_k) \geq 6$. If $O_i$ is the observed number of values in bin $i$, and $E_i$ is the number predicted by the random model, then
\[
\chi^2=\sum_i \frac{(O_i-E_i)^2}{E_i}.
\]
Thus a large value of $\chi^2$ indicates a large deviation from the proposed random model. Furthermore, it is natural to consider the normalized statistic $\chi^2/\nu$,%
where
\[
\nu = \#\{\text{Bins used in the $\chi^2$-test}\} -1.
\]
In Table \ref{tab:chi2-vs-prime} , we compare the positive $p$-valuations of our observed data with the random model $\mathbb{P} (v_p(n_k) = r) = p^{-r}$, for $r \geq 1$ (Throughout this section, $\mathbb P$ denotes the empirical proportion obtained
by sampling $k$ uniformly from the indicated set of parameters.)%
For each prime $p \geq 5$, the table gives the resulting chi-square statistic, the number of degrees of freedom, and the normalized statistic $\chi^2/\nu$. It can be seen that the prime $p = 5$ shows a strong deviation from this random model.
The exact counts are recorded in Table \ref{tab:v5-exact}. The prime $p=37$ also shows a
noticeable discrepancy, although it is substantially smaller.

\begin{table}[H]
\centering
\begin{minipage}[t]{0.42\linewidth}
\centering
\begin{tabular}{r r r r}
\toprule
$p$ & $\chi^2$ & $\nu$ & $\chi^2/\nu$\\
\midrule
5  & 76.84 & 6 & 12.81\\
7  & 1.56  & 5 & 0.31\\
11 & 8.32  & 4 & 2.08\\
13 & 6.42  & 4 & 1.60\\
17 & 1.49  & 3 & 0.50\\
19 & 4.10  & 3 & 1.37\\
23 & 3.37  & 3 & 1.12\\
29 & 3.41  & 3 & 1.14\\
31 & 7.33  & 3 & 2.44\\
37 & 14.27 & 3 & 4.76\\
41 & 2.60  & 2 & 1.30\\
43 & 2.41  & 2 & 1.20\\
47 & 1.38  & 2 & 0.69\\
\bottomrule
\end{tabular}
\subcaption{Comparison of $v_p(n_k)$ to $\mathbb{P}(v_p=r)=p^{-r}$,
primes $5\le p \le47$. }
\label{tab:chi2-vs-prime}
\end{minipage}
\hfill
\begin{minipage}[t]{0.54\linewidth}
\centering
\begin{tabular}{c r r c}
\toprule
$v_5(n_k)= r$ & count & expected & ratio\\
\midrule
0 & $186{,}238$ & $187{,}499.3$ & $0.993$ \\
1 &  $50{,}398$ &  $49{,}999.8$ & $1.008$ \\
2 &  $10{,}651$ &  $10{,}000.0$ & $1.065$ \\
3 &   $2{,}141$ &   $2{,}000.0$ & $1.071$ \\
4 &         $471$ &         $400.0$ & $1.178$ \\
5 &          $78$ &          $80.0$ & $0.975$ \\
6 &          $19$ &          $16.0$ & $1.188$ \\
7 &           $3$ &           $3.2$ & $0.938$ \\
\bottomrule
\end{tabular}
\subcaption{ Counts with $v_5(n_k)=r$ vs the expected $M\cdot 5^{-r}$, where $M:=249{,}999$  (the expected total at the $r=0$ row is $M\cdot\tfrac34$).}
\label{tab:v5-exact}
\end{minipage}
\caption{The law $\mathbb{P}(v_p=r)=p^{-r}$ ($r\ge1$) compared to our observed data. Left: reduced $\chi^2$ across primes, only $p=5$ deviates strongly. Right: exact counts $v_5(n_k)=r$.}
\label{tab:chi2-and-v5-exact}
\end{table}


\subsection{$3$-valuation observations}
The $3$-adic behavior of $n_k$ does not seem to follow the same pattern as the one described in Section~\ref{subsection:big_p-adic}. We record several observations, including parity laws for the $3$-valuations.

\subsubsection*{Globally higher $3$-divisibility.}

Divisibility by $3$ appears to occur considerably more frequently than the corresponding heuristic would suggest with $v_3(n_k)\geq 1$ about $73\%$ of the time; see  Table~\ref{tab:v3_statistics_cumulative} and Figure~\ref{fig:histogram_v3_0_250K}. More specifically, the proportions of $k$ for which $v_3(n_k)=r$ with $r \geq 2$ appear to be elevated (by at least a factor of $2$) compared with the baseline prediction $3^{-r}$. On the other hand, the proportion with $v_3(n_k)=1$ appears to be slightly smaller than this baseline: for instance, our data give approximately $0.322$, compared with $1/3 = 0.333\ldots$ expected from the model for $p\geq 5$. This excess is not an artifact of
the congruence phenomena discussed below.

\subsubsection*{$3$-Divisibility among congruence classes.}

The $3$-valuation is further structured by congruences. The proportion of $k$ with
$v_3(n_k)\geq r$ is systematically higher for $k\equiv 0\pmod 3$; see
Table~\ref{tab:v3_geq_mod3_statistics_congruence_cumulative}. Refining modulo $27$, we
find a stronger statement: for $k\equiv 0,\pm 5\pmod{27}$
one has $v_3(n_k)\geq 1$ for all but a few exceptional $k$ in the dataset.


\medskip
\noindent\emph{Bad primes of $E_k$ and $3$-divisibility.} Another observation is that there seems to be a correlation between the parity of the cardinality of the set
\[
S_k := \{ \text{bad primes of the form }p \equiv 1 \, \pmod 3 \text{ for the curve }E_k \}
\]
and the $3$-valuation of $n_k$, see Table \ref{tab:contingency_parity_S_T}. The table suggests that it is more likely for $n_k$ to be divisible by $3$ when the parity of $\#S_k$ is even: We have $\mathbb{P}(v_3(n_k) \geq 1 \, | \, \# S_k \text{ is even}) = 0.79$, compared to $\mathbb{P}(v_3(n_k) \geq 1 \, | \, \# S_k \text{ is odd}) = 0.67 $.%
If we compare this to the analogous set
\[
T_k := \{ \text{bad primes of the form }p \equiv 2 \, \pmod 3 \text{ for the curve }E_k \},
\]
this effect disappears. We have $\mathbb{P}(v_3(n_k) \geq 1 \, | \, \# T_k \text{ is even}) = 0.73$, compared to $\mathbb{P}(v_3(n_k) \geq 1 \, | \, \# T_k \text{ is odd}) = 0.73$.
This phenomenon seems to be a consequence of a sharper law on a subfamily, which we will now describe.

\begin{table}[ht]
  \centering
  \small
  \setlength{\tabcolsep}{4pt}

  \begin{tabular}{lrrr}
    \toprule
    & $v_3(n_k) \leq 0$ & $v_3(n_k)\geq 1$ & total \\
    \midrule
    $\#S_k$ even
    & \countpct{$26{,}170$}{$20.9$}
    & \countpct{$98{,}896$}{$79.1$}
    & $125{,}066$ \\
    $\#S_k$ odd
    & \countpct{$40{,}811$}{$32.7$}
    & \countpct{$84{,}122$}{$67.3$}
    & $124{,}933$ \\
    \midrule
    total
    & \countpct{$66{,}981$}{$26.8$}
    & \countpct{$183{,}018$}{$73.2$}
    & $249{,}999$ \\
    \bottomrule
  \end{tabular}
\qquad
  \begin{tabular}{lrrr}
    \toprule
    & $v_3(n_k) \leq 0$ & $v_3(n_k)\geq 1$ & total \\
    \midrule
    $\#T_k$ even
    & \countpct{$33{,}406$}{$26.7$}
    & \countpct{$91{,}600$}{$73.3$}
    & $125{,}006$ \\
    $\#T_k$ odd
    & \countpct{$33{,}575$}{$26.9$}
    & \countpct{$91{,}418$}{$73.1$}
    & $124{,}993$ \\
    \midrule
    total
    & \countpct{$66{,}981$}{$26.8$}
    & \countpct{$183{,}018$}{$73.2$}
    & $249{,}999$ \\
    \bottomrule
  \end{tabular}

  \caption{Contingency tables comparing the divisibility of $n_k$ by $3$ with the parities of $\#S_k$ and $\#T_k$. Percentages are according to the parity of
$\#S_k$ and $\#T_k$.}
  \label{tab:contingency_parity_S_T}
\end{table}


\subsubsection{A parity law on a subfamily}
\label{ssec:partitylaw}

On the set
\begin{equation}
\label{def:setA}
A := \{ k \, \colon \, 3\nmid k, \text{ and } k(k^2-16)\text{ squarefree} \},
\end{equation}
the above-mentioned correlations sharpen into two laws, valid (with the exception of $k = 1$):
\begin{equation}\label{eq:thelawforA}
v_3(n_k)=0 \Longrightarrow \#S_k \text{ is odd},
\qquad
v_3(n_k)=1 \Longrightarrow \#S_k \text{ is even}.
\end{equation}
see Tables \ref{tab:parity-law} and \ref{tab:law-contingency}. Here it can be seen that a similar law does not seem to appear for $v_3(n_k) \geq 2$.


\begin{table}[ht]
  \centering
  \begin{tabular}{rrrrr}
    \toprule
    $v_3(n_k)$ & count & $\#S_k$ even & $\#S_k$ odd & $\mathbb{P}(\#S_k\ \mathrm{odd})$ \\
    \midrule
    0 & $13{,}080$ & $1$ & $13{,}079$ & $1.000$ \\
    1 & $13{,}233$ & $13{,}233$ & $0$ & $0.000$ \\
    2 & $8{,}950$ & $4{,}438$ & $4{,}512$ & $0.504$ \\
    3 & $4{,}161$ & $2{,}055$ & $2{,}106$ & $0.506$ \\
    4 & $1{,}560$ & $780$ & $780$ & $0.500$ \\
    5 & $552$ & $266$ & $286$ & $0.518$ \\
    6 & $194$ & $95$ & $99$ & $0.510$ \\
    7 & $64$ & $31$ & $33$ & $0.516$ \\
    8 & $22$ & $10$ & $12$ & $0.545$ \\
    9 & $9$ & $5$ & $4$ & $0.444$ \\
    10 & $3$ & $1$ & $2$ & $0.667$ \\
    11 & $1$ & $0$ & $1$ & $1.000$ \\
    12 & $1$ & $1$ & $0$ & $0.000$ \\
    \bottomrule
  \end{tabular}
    \caption{Parity of $\#S_k$ as a function of $v_3(n_k)$, for $k \in A$. The law $v_3(n_k)=0\Rightarrow\#S_k$ odd and $v_3(n_k)=1\Rightarrow\#S_k$ even holds with a single exception; for $v_3(n_k)\ge2$ the parity is unconstrained.}
        \label{tab:parity-law}
\end{table}

\begin{table}[ht]
  \centering
  \begin{tabular}{lrrr}
    \toprule
               & $v_3(n_k) = 0 $ & $v_3(n_k) \geq 1 $ & total \\
    \midrule
    $\#S_k$ even  & $1$ & $20{,}915$ & $20{,}916$ \\
    $\#S_k$ odd   & $13{,}079$ & $7{,}835$ & $20{,}914$ \\
    \midrule
    total       & $13{,}080$ & $28{,}750$ & $41{,}830$ \\
    \bottomrule
  \end{tabular}
    \caption{Contingency of $\mathrm{parity}(\#S_k)$ against $v_3(n_k) \geq 1$ for $k \in  A$.}
  \label{tab:law-contingency}

\end{table}


It is apparent that law \eqref{eq:thelawforA} is a special case of a more general law that holds on a larger set. Define
\[
\alpha(k) := \#S_k + v_3(k) + \mathbbm{1}_{8 \mid k}.
\]
The two laws are
\begin{subequations}\label{eq:law-general}
\begin{align}
\text{(odd)}&\qquad v_3(n_k)=0 \ \Longrightarrow\ \alpha(k)\ \text{is odd},
   \label{eq:law-fwd}\\
\text{(even)}&\qquad v_3(n_k)=1 \ \Longrightarrow\ \alpha(k)\ \text{is even}.
   \label{eq:law-conv}
\end{align}
\end{subequations}
On $A$ we have $k^2-16$ is squarefree, $v_3(k) = 0$ and $8 \nmid k$, so \eqref{eq:law-general} reduces to \eqref{eq:thelawforA}. We now investigate several enlargements of the set $A$ for which the laws \eqref{eq:law-general} hold, see Table~\ref{tab:sets} and Figure~\ref{fig:placeholder}. The results in Table~\ref{tab:sets} suggest that, among the sets considered there, the largest set on which the  odd law~\eqref{eq:law-fwd} holds is
\[
C' := \{p^2 \nmid k^2 - 16 \text{ for all }p \equiv 3 \pmod 4 \} \setminus \{k \equiv 4 \pmod 8 \}.
\]
On this set, however, the even law \eqref{eq:law-conv} has $1{,}746$ exceptions. Of these, $1{,}597$ are divisible by $3^3$, and hence satisfy $v_3(n_k) \geq 1$ by the earlier remarks. For the remaining $149$, the integer $k(k^2-16)$ is divisible by one of $\{7^3, 13^3, 19^3, 37^3, 61^3 \}$, all cubes of primes congruent to $1 \pmod 3$. This suggests that both laws \eqref{eq:law-general} may hold on the smaller set
\[
C'' = \{k \in C' \, \colon\, p^3 \mid k(k^2-16) \Rightarrow p \equiv 2 \pmod 3 \}.
\]
Finally, it is worth noting that the $k$ such that $k \equiv 0 \pmod 27$ or $k \equiv \pm 5 \pmod 27$ lie outside the set $C''$:  the former because $3^3 \mid k$ and the latter because $9 \mid k^2-16$ obstructs squarefreeness at the prime  $3 \equiv 3 \pmod 4$. So their elevated $3$-valuation is not explained by the law on $C''$.
\begin{table}[h]\centering
\renewcommand{\arraystretch}{1.25}
\begin{tabular}{@{}lll cc@{}}
\toprule
set & definition & density & \eqref{eq:law-fwd} exc. & \eqref{eq:law-conv} exc.\\
\midrule
$A$      & $3\nmid k,\ k(k^2-16)$ squarefree      & $0.167$ & $0$ & $0$\\
$A'$     & $3\nmid k,\ k^2-16$ squarefree         & $0.184$ & $0$ & $46$\\
$A''$    & $k^2-16$ squarefree                    & $0.323$ & $0$ & $913$\\
$B$ & odd part of $k^2 - 16$ squarefree & $0.645$ & $3{,}076$ & $4{,}753$\\
$B'$ & $B \setminus \{ k \equiv 4 \pmod 8\}$ & $0.565$ & $0$ & $1{,}538$\\
$B''$ & $B' \cap \{p^3 \mid k\Rightarrow p\equiv 2 \pmod 3 \}$ & $0.536$ & 0 & 0  \\
$C$ & $p^2 \nmid k^2 - 16$ for all $p \equiv 3 \pmod 4$ & $0.721$ & $3{,}423$ &  $5{,}330$ \\
$C'$ & $C \setminus \{ k \equiv 4 \pmod 8 \}$ & $0.631$ & $0$ &  $1{,}746$ \\
$C''$ & $C' \cap  \{p^3 \mid k(k^2 - 16) \Rightarrow p\equiv 2 \pmod 3 \}$ & $0.598$ & $0$ & $0$ \\
 & all values $k$ & 1.000 & $11{,}342$ & $17{,}588$ \\
\bottomrule
\end{tabular}
\caption{The parity laws \eqref{eq:law-general} on different enlargements of $A$, computed over the range $5 \leq k\leq 250{,}000$. Here ``density'' is the proportion of integers $k$ that lie in the set, and ``exc.'' is the number of exceptions of the corresponding law.}
\label{tab:sets}
\end{table}
\begin{figure}[h]
  \centering
  \newcommand{\SCALE}{20}%
  \begin{tikzpicture}[
      >=stealth,
      every node/.style={font=\normalsize},
      inc/.style={thin,->},
      dot/.style={circle,fill=black,inner sep=1.3pt}
  ]
  \foreach \d in {0.2,0.3,0.4,0.5,0.6,0.7}{
     \node[gray,font=\small,left] at (-1.0,\d*\SCALE) {\d};
  }
 \draw[->,gray] (-1.0,0.15*\SCALE) -- (-1.0,0.76*\SCALE) node[above,black,font=\small] {density};
  \node[dot,label=below:{$A$}]   (A)   at (0,    0.167*\SCALE) {};
  \node[dot,label=left:{$A'$}]   (Ap)  at (0,    0.184*\SCALE) {};
  \node[dot,label=left:{$A''$}]  (App) at (0,    0.323*\SCALE) {};
  \node[dot,label={[label distance=2pt]-30:{$B''$}}]  (Bpp) at (3.2,  0.536*\SCALE) {};
  \node[dot,label=left:{$B'$}]   (Bp)  at (3.2,  0.565*\SCALE) {};
  \node[dot,label=left:{$B$}]    (B)   at (3.2,  0.645*\SCALE) {};
  \node[dot,label=right:{$C''$}] (Cpp) at (6.4,  0.598*\SCALE) {};
  \node[dot,label=right:{$C'$}]  (Cp)  at (6.4,  0.631*\SCALE) {};
  \node[dot,label=right:{$C$}]   (C)   at (6.4,  0.721*\SCALE) {};
  \draw[inc] (A) -- (Ap);   \draw[inc] (Ap) -- (App);  \draw[inc] (A) -- (Bpp);
  \draw[inc] (App) -- (Bp); \draw[inc] (Bpp) -- (Bp);  \draw[inc] (Bp) -- (B);
  \draw[inc] (Cpp) -- (Cp); \draw[inc] (Cp) -- (C);    \draw[inc] (Bpp) -- (Cpp);
  \draw[inc] (Bp) -- (Cp);  \draw[inc] (B) -- (C);
  \end{tikzpicture}
  \caption{The sets of Table~\ref{tab:sets}, placed at height equal to their density.
  An arrow points from a set to a larger set containing it; since a subset is never
  denser than its superset, every arrow points upward.}
  \label{fig:placeholder}
\end{figure}


\subsection{$2$-adic observations}

Whenever $n_k$ is an integer it appears to be even, the only exceptions being
$k=1,2$ (where $n_k=1$). More is true: among integer values of $n_k$, the
inequality $v_2(n_k)\geq 2$ holds for all $k$ in the dataset outside the finite set
\[
k \in \{1,2,6,7,9,20,23,28,36,48,96,112,432\}.
\]
The distribution of $v_2(n_k)$ is recorded in
Table~\ref{tab:v2_statistics_cumulative} and
Figure~\ref{fig:histogram_v2_0_250K}.


\medskip
\noindent\emph{Dependence on congruence classes.} We next consider, for a fixed odd prime $q$, the proportion of $k$ in a residue
class $k\equiv a \pmod q$ with $v_2(n_k)\geq r$. For small $r$ this proportion
is close to $1$ in every class, so the structure only becomes visible for
larger $r$. Empirically, it is largest when $a\equiv\pm4\pmod q$ and second
largest when $a\equiv 0\pmod q$; see for example
Table~\ref{tab:v2_geq_mod7_statistics_congruence_cumulative}, where the
frequency of $v_2(n_k)\geq r$ is highest for $k\equiv\pm4\pmod 7$, closely
followed by $k\equiv0\pmod 7$. These are exactly the classes of bad reduction:
since $\Delta_k=2^{-24}k^{14}(k^2-16)$, we have
$\Delta_k\equiv0\pmod q$ if and only if $k\equiv 0,\pm4\pmod q$.

\medskip
\noindent\emph{Empirical model.} From the earlier observations, the primes dividing $k^2-16$ (the case
$k\equiv\pm4\pmod q$) have more influence than those dividing $k$ (the case
$k\equiv0\pmod q$), and empirically the appropriate weights are $2$ and $1$
respectively. We therefore assign to each odd prime $q$ the weight
\[
  W_q(k) := \begin{cases}
0 & q\nmid k(k^2-16),\\
1 & q\mid k,\\
2 & q\mid k^2-16,
\end{cases}\in\{0,1,2\},
\]
well defined since an odd prime divides at most one of $k,\,k-4,\,k+4$. Writing $\omega_{\mathrm{odd}}(n)$ for the number of
distinct odd prime divisors of $n$, we have
\[\sum_{q\ \mathrm{odd}} W_q(k)=\omega_{\mathrm{odd}}(k)+2\,\omega_{\mathrm{odd}}(k^2-16)=\omega_{\mathrm{odd}}(k)+2\,(\omega_{\mathrm{odd}}(k-4)+\omega_{\mathrm{odd}}(k+4)),\]
and find a moderately strong positive correlation (Pearson $r=0.69$) between this
quantity and $v_2(n_k)$. Since this sum involves only the odd primes, it appears to be natural to correct it by a term $c(k)$ that incorporates the local behavior of $E_k$ at the prime $2$. Let
\begin{equation}\label{eq:c-defn}
  c(k) =
  \begin{cases}
    4, & E_k \text{ semistable at } 2\ \ (v_2(N_k)\in\{0,1\}),\ \text{i.e. } k \text{ odd or } 16\mid k,\\
    2, & E_k \text{ additive at } 2\ \ (v_2(N_k)\in\{3,4\}),\ \text{otherwise.}
  \end{cases}
\end{equation}
We then set
\begin{equation}\label{eq:hatvdef}
  \hat v(k) := \sum_{q\ \mathrm{odd}} W_q(k) - c(k)
  = \omega_{\mathrm{odd}}(k) + 2\,\omega_{\mathrm{odd}}(k^2-16) - c(k).
\end{equation}
Then
$v_2(n_k)$ and $\hat v(k)$ have Pearson correlation $r=0.77$
(Figure~\ref{fig:correlationv2}), improving on the value $r=0.69$ obtained
without the correction $c(k)$. In fact, $\hat{v}(k)$ seems to underestimate $v_2(n_k)$ for
all but finitely many $k$.
\begin{figure}[ht]
  \centering
\includegraphics[width=0.85\textwidth]{correlation_v2.pdf}
  \caption{ $v_2(n_k)$ vs
    $\hat{v}(k)$ for $k$ with $v_2(n_k) \geq 0$}
  \label{fig:correlationv2}
\end{figure}
This leads us to the following conjecture.
\begin{conj}
\label{conj:low_bound}
For all $k \neq 4, 8$, we have
    \[\hat{v}(k) \leq v_2(n_k).\]
\end{conj}


Moreover, the excess $G(k) := v_2(n_k) - \hat{v}(k)$ has mean $2.05 $ and the data suggest that $G(k)=0$ may occur infinitely often, making the conjecture sharp in that sense.
This excess seems to spike whenever $v_2(k) = 3$, and then drops down below $2$ when $v_2(k) \geq 4$ (see Table \ref{tab:v2-left}). This spike comes from $k \equiv \pm 8 \pmod{64}$, where the mean excess is $3.02$ (see Table \ref{tab:v2-right}). In fact it grows further as $k \equiv \pm 8 \pmod{2^m}$ for increasing $m \geq 6$, plateauing at an excess of $3.5$ (see Table \ref{tab:v2-bottom}).
\begin{table}[h]


\centering
\begin{minipage}[t]{0.52\linewidth}
\centering
\begin{tabular}{l r r c}
\toprule
class & count & mean $v_2(n_k)$ & mean $G(k)$\\
\midrule
$k$ odd        & 125{,}000 & 9.56  & $+2.05$\\
$v_2(k)=1$     & 62{,}500  & 11.28 & $+2.04$\\
$v_2(k)=2$     & 31{,}249  & 10.39 & $+2.00$\\
$v_2(k)=3$     & 15{,}625  & 11.29 & $+2.41$\\
$v_2(k)=4$     & 7{,}813   & 8.67  & $+1.86$\\
$v_2(k)=5$     & 3{,}906   & 8.65  & $+1.91$\\
$v_2(k)=6$     & 1{,}953   & 8.55  & $+1.90$\\
$v_2(k)=7$     & 977    & 8.52  & $+1.93$\\
$v_2(k)\ge 8$  & 976    & 8.33  & $+1.95$\\
\midrule
all            & 249{,}999 & 10.15 & $+2.05$\\
\bottomrule
\end{tabular}
\subcaption{$v_2(n_k)$ by the $2$-adic class of $k$.}
\label{tab:v2-left}
\end{minipage}%
\hspace{0.0015\linewidth}%
\begin{minipage}[t]{0.42\linewidth}
\centering
\begin{tabular}{c r c c}
\toprule
$k\pmod 64$ & count & mean $v_2(n_k)$ & mean $G(k)$\\
\midrule
$8$  & $3{,}907$ & $11.96$ & $+3.09$\\
$24$ & $3{,}906$ & $10.69$ & $+1.81$\\
$40$ & $3{,}906$ & $10.70$ & $+1.82$\\
$56$ & $3{,}906$ & $11.83$ & $+2.94$\\
\midrule
all $v_2(k)=3$ & $15{,}625$ & $11.29$ & $+2.41$\\
\bottomrule
\end{tabular}
\subcaption{The class $v_2(k)=3$, refined by $k \pmod {64}$.}
\label{tab:v2-right}
\end{minipage}

\vspace{1em}

\begin{minipage}[t]{0.62\linewidth}
\centering
\begin{tabular}{c r r r r}
\toprule
$M$ & count & mean $v_2(n_k)$ & mean $G(k)$ & $2+s(k)$\\
\midrule
$5$        & $7{,}812$ & 10.69 & $+1.81$ & $2.00$\\
$6$        & $3{,}906$ & 11.53 & $+2.65$ & $2.75$\\
$7$        & $1{,}954$ & 12.19 & $+3.31$ & $3.13$\\
$8$        & $976$  & 12.30 & $+3.43$ & $3.31$\\
$9$        & $488$  & 12.39 & $+3.50$ & $3.41$\\
$\ge 10$   & $489$  & 12.37 & $+3.47$ & $\to 3.5$\\
\bottomrule
\end{tabular}
\subcaption{Excess by exact $2$-adic depth $M:=\max\{v_2(k-8),v_2(k+8)\}$,
compared with the model mean $2+s(k)$.}
\label{tab:v2-bottom}
\end{minipage}

\caption{$2$-adic valuation $v_2(n_k)$, with excess
$G(k)=v_2(n_k)-\hat v(k)$. Left: grouped by the $2$-adic class of $k$.
Right: the exceptional class $v_2(k)=3$, split further modulo $64$. Bottom: the exceptional classes by $2$-adic depth $M=\max\{v_2(k-8),v_2(k+8)$ \}. }
\label{tab:v2-valuation-classes}
\end{table}
A reasonable fit for the mean excess $G(k)$ seems to be
\begin{equation}\label{eq:s-defn}
\E[G(k)]\;\approx\;2+s(k),\qquad
s(k):=\tfrac{3}{2}\,\max\!\Big\{\,1-2^{\,5-\max\{v_2(k-8),\,v_2(k+8)\}},\;0\,\Big\}.
\end{equation}
Here $s(k)$ vanishes when $k$ is not $2$-adically close to $\pm 8$ and approaches $1.5$ as $k$ approaches $\pm 8$ $2$-adically.

Combining these observations, our empirical model is
\begin{equation}\label{eq:v2approx}
v_2(n_k) = \hat{v}(k) + s(k) + X_k,
\end{equation}
where the residual term $X_k$ appears to have a distribution that is largely independent of $k$. More precisely, its empirical distribution is well approximated by
\[
X_k\;\approx\;\mathrm{NB}\!\left(r=3,\;p=\tfrac{3}{5}\right),
\]
where $\mathrm{NB}$ denotes the \emph{negative binomial distribution}.
Accordingly, in this probabilistic model, we take
\[
\mathbb P(X_k=j)
 = \binom{j+2}{2}
   \left(\frac35\right)^3
   \left(\frac25\right)^j,
 \quad j\geq 0.
\]
This distribution has mean $2$ and mode $1$. A comparison between this statistical model and our data is shown in Figure~\ref{fig:histogramv2stat}.
\begin{figure}
    \centering
    \includegraphics[scale=0.5]{v2_data_vs_model.pdf}
    \caption{The observed distribution of $v_2(n_k)$ for values with $v_2(n_k)\geq 0$, compared with the prediction
    $\hat{v}(k)+s(k)+\mathrm{NB}(3,\tfrac{3}{5})$.}
    \label{fig:histogramv2stat}
\end{figure}

Since the mode of this distribution is $1$, the best deterministic predictor we found for $v_2(n_k)$ is
\begin{equation}\label{eq:P}
\hat{P}(k):=\mathrm{round}\bigl(\hat{v}(k)+s(k)+1\bigr).
\end{equation}
This predictor has Pearson correlation coefficient $0.78$ with $v_2(n_k)$, only slightly higher than the correlation obtained using $\hat{v}(k)$ alone. Moreover, the prediction $\hat{P}(k)$ is correct $24.3\%$ of the time.


\begin{rem}
Since $\hat v(k) + c(k) = \sum_{q\ \mathrm{odd}} W_q(k)$ is  expressed as a sum of local contributions over the odd primes, and the conditions
$q\mid k$ and $q\mid (k^2-16)$ hold on $1$ and $2$ residue classes modulo $q$
(with densities $1/q$ and $2/q$), each odd prime contributes
$1\cdot\tfrac1q + 2\cdot\tfrac2q = \tfrac5q$ to $\hat v$ on average. By Mertens'
theorem, and since $c(k)$ is bounded,
\[
  \frac1x\sum_{4\le k\le x}\hat v(k) = \sum_{q\le x}\frac5q + O(1)
  = 5\log\log (x) + O(1)
\]
as $x \to \infty$.  So, assuming Conjecture \ref{conj:low_bound}, the average order of $v_2(n_k)$ is at least $5 \log \log (x)$. Tables~\ref{tab:largeW-odd} and \ref{tab:largeW-even} list some of the  maximisers of $W := \omega_{\mathrm{odd}}(k)+2\,\omega_{\mathrm{odd}}(k^2-16)$ in the window
$2.5\times10^{5}<k<2\times10^{6}$. In every case $v_2(n_k) \geq  \hat{v}(k)$, with equality  at $k = 1540136$, $k=1812794$, and $k=1984814$. We also remark the large value $v_2(n_k)=27$ occurring at $k=1069814$ (though the largest value we observe is $v_2(n_k) = 31$, at the much smaller $k = 36326$).%

\end{rem}

\subsection{The regime $k=\sqrt m$}
\label{sec:sqrtm}

For $m < 0$ the parameter $k = \sqrt{m}$ is imaginary and \eqref{HG} does not apply directly; instead we use the functional equation \cite[Theorem~2.2]{LR07} (taking $k=(\sqrt{-m}-\sqrt{16-m})/4$ in their notation).
Using the notation $\mu(k):=\mu_k$, giving
\[\mu(\sqrt m)=\mu\!\left(\frac{64}{(\sqrt{-m}-\sqrt{16-m})^2}\right)-\mu(\sqrt{16-m}),\qquad m<0,\]
allowing us to compute $r_k$ also for $m<0$.

We have computed these $r_k$ for $|m| \leq 30{,}000$. We briefly compare the observations with integer values of $k$. First of all, for
\[
m \not \in \{ -768, -16, -9, -2, -1, 9, 18, 25, 27, 64, 144, 256 , 1024\},
\]
$r_k$ is the reciprocal of an integer.
For $p\ge5$, the law $\mathbb P(v_p(n_k)=r)\approx p^{-r}$ holds unchanged.
 For $p=3$, the parity law of Section \ref{ssec:partitylaw} breaks down completely (it holds
for only about half of the relevant $k$); of the accompanying congruences, only
$m\equiv25\bmod27$ survives, that is, the image of $k\equiv\pm5$. This forces once again $v_3(n_k)\ge1$.


For $p=2$, a variant of the lower bound in Conjecture \ref{conj:low_bound} holds. The
term $\omega_{\text{odd}}(k)$ appears to split into two different counts on $m$
\[
\omega^0_{\text{odd}}(m) + 2 \omega_{\text{odd}}^1(m),
\]
where $\omega^0_{\text{odd}}(m)$ is the number of odd primes $p$ dividing $m$ with
$v_p(m)$ even, and $\omega^1_{\text{odd}}(m)$ the number of odd primes with $v_p(m)$ odd.
Together with a sign-dependent shift that makes the bound sharp, the generalization of
$\hat{v}(k)$ becomes
\[
\hat{v}(m) = \omega^0_{\text{odd}}(m) + 2 \omega^1_{\text{odd}}(m)
             + 2 \omega_{\text{odd}}(m - 16) - c(m) + \varepsilon(m),
\]
where
\[
c(m) =
  \begin{cases}
    4, & v_2(m) = 0 \text{ or } v_2(m) \geq 8, \\
    2, & 1 \leq v_2(m) \leq 7,
  \end{cases}
\qquad
\varepsilon(m) =
  \begin{cases}
    0, & |m|\ \text{a perfect square},\\
    1, & m<0\ \text{nonsquare},\\
    2, & m>0\ \text{nonsquare}.
  \end{cases}
\]
Then $\hat{v}(m) \leq v_2(n_k)$ appears to hold for all $m$ with integral $n_k$, the sole
exception being $m=12$. Equality   is attained for about $17\%$ of these $m$; the
residual $v_2(n_k) - \hat{v}(m)$ is then non-negative, of the same negative-binomial
flavor as the excess $X_k$ in the integer case. We further remark
that all of the behavior described above also holds when $m$ is restricted to a single sign.


\subsection{A conditional proof of Conjecture \ref{conj:low_bound} for special families of $k$}
\label{ssec:lowboundproof}
We now use the work of Dummigan, Golyshev, de Jeu and Kerr
\cite{DGJK} to verify Conjecture~\ref{conj:low_bound} conditionally
on the Bloch--Kato conjecture and the additional hypotheses appearing
in their work.  Let
\[
\widetilde E_u\colon y^2=x(x+1)(x+u^2).
\]
The curve $\widetilde E_u$ is $2$-isogenous to $E_{4u}$; hence, when
$k=4u$, the two curves have the same $L$-function and conductor.
\begin{prop}
\label{deJeuprop}
        Let $k = 4 u$, where $u \equiv 4 \pmod 8$ and $u(u^2-1)$ is squarefree at the odd primes; assume in addition that $u+1$ has a prime divisor $\equiv 3 \pmod 4$.
            Assume moreover,
as in \cite[Thm.~3]{DGJK}, that $K_2^T(\tilde E_u)_{\mathrm{int}}\otimes_{\mathbb Z}\mathbb Q$
is one-dimensional, that $r_k$ is a rational number and that $H^1_f(\mathbb Q,\tilde E_u[2^\infty](-1))$ is finite;
and assume the Bloch--Kato conjecture for $\tilde E_u$.

Then Conjecture~\ref{conj:low_bound}, i.e. $\hat v(k)\le v_2(n_k)$, is true for these $k$.
    \end{prop}

\begin{proof}
By the hypotheses on $u$, the hypothesis of \cite[Theorem 1]{DGJK} are satisfied and $m_u = 1$, so that $\iota_{u,2} = 1$. In the notation of \cite{DGJK}, we have $F(u)=\mu_k$ and
\[
q_u=\frac{L(\widetilde E_u,2)}{(2\pi i)^2F(u)}.
\]
By the functional equation and the fact that $\widetilde E_u$ is
isogenous to $E_k$, one has
\[
n_k=-\omega_kN_kq_u.
\]
Under the present hypotheses, the conductor is
\[
N_k=\frac{u(u^2-1)}{4},
\]
which is odd.  Consequently,
\[
v_2(n_k)=v_2(q_u).
\]
    On the one hand, since $\iota_{u,2} = 1$, \cite[Equation (1.3)]{DGJK} reads
    \begin{equation}
    \label{v2_deJeu}
             v_2(n_k) + 2 = 2 \omega_1(u) + \omega_3(u) + \omega(u^2 -1) + s_u,
    \end{equation}
    where
    \[
    s_u = v_2 \left( \#H_f^{1}(\Q, \tilde{E}_{u}[2^\infty](-1))\right),
    \]
     $\omega$ is the number of distinct prime factors, and $\omega_j$ the number of distinct prime factors $\equiv j \pmod 4$. On the other hand, $k = 4u$ implies $\omega_{\text{odd}}(k) = \omega_{\text{odd}}(u) = \omega_1(u) + \omega_3(u)$ and similarly, $\omega_{\text{odd}}(k^2 - 16) = \omega(u^2-1)$. Since  $16 | k$, we have $c(k) = 4$.
Hence
\begin{equation}
\label{hat_in_u}
\hat{v}(k) = \omega_1(u) + \omega_3(u) + 2 \omega(u^2 - 1) - 4 .
\end{equation}
Combining the two Equations \eqref{v2_deJeu} and \eqref{hat_in_u} gives
\[
v_2(n_k) - \hat{v}(k) = s_u - (\omega(u^2 - 1)-\omega_1(u)- 2).
\]
Thus, in order to prove our conjecture for these $k$, it suffices to prove the inequality
\[
s_u \geq \omega(u^2 - 1)-\omega_1(u)- 2.
\]
By elementary group theory we have that
\[
s_u \geq s'_u :=v_2(\#H_f^{1}(\Q, \tilde{E}_{u}[2^\infty](-1))[2]).
\]
We now use the explicit description of this $2$-torsion subgroup given in \cite[Theorem 2]{DGJK}. Let $S$ be the set of prime divisors of $u^2-1$, and let $S'$ be the set of prime divisors of $u$ that are congruent to $1 \pmod 4$. Restricting the parametrization in \cite[Theorem~2]{DGJK} to pairs of the form $(D,1)$ as discussed in \cite[Remark~7.9]{DGJK}, we obtain a subspace consisting of the squarefree divisors $D$ of $u^2-1$ such that
\[
D \text{ is a square modulo every prime in } S'
\qquad\text{and}\qquad
D\equiv 1\pmod 8.
\]
Writing $D=\prod_{p\in S}p^{x_p}$ with $x_p\in \F_2$, the above two conditions give a linear system in the $\omega(u^2 - 1)$ variables $x_p$. Indeed, if $\left( \frac{p}{q} \right) = (-1)^{\alpha_{p, q}}$ for $\alpha_{p, q} \in \mathbb{F}_2$ for each prime $q$ dividing $u$ such that $q \equiv 1 \pmod 4$, we get a system of $\omega_1(u)$ equations $\sum_{p} x_p \alpha_{p,q} = 0$ in $\omega(u^2 - 1)$ variables for the first condition.  The second condition can be expressed as $\left(\frac{-1}{D}\right)=\left(\frac{2}{D}\right)=1$, and this gives at most two extra equations, $\sum_p x_p\beta_{p,-1}=\sum_p x_p\beta_{p,2}=0$, where $\left(\frac{-1}{p}\right)=(-1)^{\beta_{p,-1}}$ and  $\left(\frac{2}{p}\right)=(-1)^{\beta_{p,2}}$.
Thus, the system has $\omega(u^2 - 1)$ variables and at most $\omega_1(u) + 2$ equations. Its solution space therefore has dimension at least   $\omega(u^2-1)-\omega_1(u)-2$. Consequently,
\[
s_u \geq \omega(u^2-1)-\omega_1(u)-2,
\]
and we recover  that statement of Conjecture \ref{conj:low_bound} under these particular conditions.
\end{proof}
\begin{rem}[The Bloch--Kato formula at an arbitrary prime]
\label{rem:general-BK}
The formulation in \cite[\S12]{DGJK} is not restricted to the prime
$2$.
Whenever the hypotheses in Assumption~12.1 of \cite{DGJK} hold, the
Bloch--Kato conjecture predicts
\[
\begin{aligned}
v_p(n_{k})
={}&v_p(N_{k})
+\sum_{q\leq\infty}
v_p\!\left(
\operatorname{Tam}_{q,\omega}^{0}(H^1_{\mathrm{\acute et}}
\bigl(\widetilde E_{u,\overline{\Q}},\Z_p\bigr)(2))
\right)\\
&+
v_p\!\left(
\#H_f^1\bigl(\Q,\widetilde E_u[p^\infty](-1)\bigr)
\right)\\
&-
v_p\!\left(
\#H^0\bigl(\Q,\widetilde E_u[p^\infty](1)\bigr)
\right)
-
v_p\!\left(
\#H^0\bigl(\Q,\widetilde E_u[p^\infty](-1)\bigr)
\right)\\
&-v_p(\iota_{u,p}),
\end{aligned}
\]
where $k=4u$ and $\iota_{u,p}$ is the regulator index defined in
\cite[\S12]{DGJK}. Thus, conjecturally, the valuation of $n_k$
decomposes into conductor and local Tamagawa contributions, a
Tate-twisted Bloch--Kato Selmer contribution, global invariant
terms, and a regulator index.


At present, however, this does not give a practically computable
formula at a general odd prime: neither the full $p$-primary
Bloch--Kato Selmer group nor the regulator index
$\iota_{u,p}$ is known in general.  For example, when $u=4$,
the identity $n_{16}=1/11$ is interpreted in \cite{DGJK} as
predicting that
\[
11\mid\iota_{4,11}.
\]
\end{rem}

\section{The model} \label{sec:model}

We train a transformer-based model \cite{transformer} using \textsc{Axolver}, a system based on Int2Int \cite{Axolver}, on our dataset. The transformer-based model has 4 encoder layers and 4 decoder layers, 8 attention heads and an embedding dimension of 256 on batches of 32 examples, using the Adam optimizer \cite{adam} with a constant learning rate of $10^{-4}$. All other parameters were left at their default values unless otherwise indicated.%

Our input is an integer $k$, encoded as a sequence of digits. The target output is a rational number $r_k$, written as a pair of integers, each again encoded as a sequence of digits. Thus, the task is to learn
\[
k \mapsto r_k
\]
from our dataset.

We randomly split our set into \emph{training}, \emph{validation} and \emph{test sets},  of sizes $243{,}999$, $3{,}000$ and $3{,}000$, respectively. The training set is used to train the neural network, while the test set is used to monitor how well the model learns over the course of training. Since the hyperparameter tuning did not yield any improvements, we did not use the validation set in our analysis. We measure the accuracy on this set by the \emph{greedy accuracy}, which is the percentage of times the model guesses correctly. We evaluate the model performance at the end of every epoch which we set to be equal to 1,000 steps.

We attempt to learn $r_k$ in three types of tasks: learning $r_k$ directly, learning the size of $|r_k|$, and learning the $p$-adic valuations $v_p(n_k)$. The tasks involving the learning of $r_k$ and of the size of $|r_k|$ are performed on base $10$ (as a default), but the tasks involving the learning of $v_p(n_k)$ are performed on base $p$, since this allows the model to detect multiplicity by $p$ and the value of $v_p$ more easily than in other bases. For the base $10$ tasks, we use \texttt{max\_len} $=256$ and \texttt{max\_output\_len} $=64$. For the valuation tasks, which are performed on bases 2, 3, 5, and 7, we use \texttt{max\_len} = \texttt{max\_output\_len} $=512$ as the tokenized data sets are too large to be handled with smaller parameters.


We run the experiments with different features (computed with \textsc{PARI/GP} \cite{PARI}) consisting of various arithmetic invariants of the elliptic curve $E_k$. These are as follows.


\begin{itemize}
\item $\omega_k$, the root number, which is the sign of the functional equation
of $L(E_k,s)$. As explained in the introduction, this value is expected to codify the sign of $r_k$.


\item $N_k$ and $\text{factor}(N_k)$, the conductor of $E_k$ and its factorization. Our goal is to measure whether information on  prime reduction can help the model determine arithmetic properties of $r_k$. We are also interested in seeing if the model can learn the size $|r_k|$ from $N_k$, in view of the estimate \eqref{orderofnk}.


\item $\Delta_k$ and $\text{factor}(\Delta_k)$, the  discriminant of $E_k$ given by the model \eqref{eq:Weierstrass}, and its factorization. Its formula is $\Delta_k=2^{-24}k^{14}(k^2-16)$. Once again, we are interested in studying the connection of prime reduction and the arithmetic properties of $r_k$. Moreover, as equations \eqref{eq:hatvdef} and \eqref{eq:v2approx} indicate, the conjectural estimate for $v_2(n_k)$ seems to involve a count on the primes of bad reduction, with odd primes counted with different weight according to the type of reduction, which is better encoded by the exponents codified in the factorization of the discriminant.

\item $a_p$, the trace of Frobenius for $p=2, 3, 5,
7, 11, 13, 17, 19, 23, 29, 31, 37$, given as a vector. We wish to observe if giving the first traces of Frobenius can help the model learn $r_k$, given that they codify the first coefficients of the $L$-function.

\item The rank of $E_k$. This was computed with the \texttt{ellrank} command in \textsc{PARI/GP} \cite{PARI} and was therefore   represented as a vector $[\mathrm{rank}_{\min{}},\mathrm{rank}_{\max{}}]$%
While the rank is supposed to codify information about $L(E_k,s)$ at $s=1$ and not $s=2$ (via the Birch and Swinnerton-Dyer conjecture), this is still a natural quantity to consider due to its importance.

\item $c_2$, the Tamagawa number at $p=2$. We consider this quantity in the learning of $v_2(n_k)$, due to its role in codifying information on the reduction of $E_k$ at $2$, which is very different from the situation of the other primes.


\end{itemize}


\section{Machine learning results}\label{sec:machine}
In this section we analyze the output of our neural network, and compare it to the results of the previous section \ref{sec:statistics}.

\subsection{The size of $|n_k|$} The initial experiments intended for the model to learn $r_k$ given various features discussed in Section \ref{sec:model} (all written in base 10).  However, this approach proved to be too ambitious, as the experiments yielded validation accuracy essentially 0. For these experiments, we noticed that  the model only seems to learn the size of $r_k$ (or equivalently, the size of $n_k$). In the discussion that follows, we will refer to {\em learning the size of $|n_k|$} to the experiment where the model output $s_k$ is considered correct if $s_k$ has the same sign as $r_k$ and if $\frac{|r_k|}{2}< |s_k|< |2r_k|$, that is, the prediction of the {\em factor-two accuracy for $r_k$}.

The first observation is that including the root number $\omega_k$ as a feature doubles the accuracy, see Figure~\ref{fig:signvsnosign}, which suggests that the model does not learn anything about the sign of $r_k$ unless it is provided with said sign as a feature. Because of this, we conducted the experiments almost always including $\omega_k$ as a feature.

\begin{figure}
    \centering
    \includegraphics[scale = 0.5]{rk_two_experiments__3_.pdf}
    \caption{The model's accuracy on the size of $n_k$ with $\omega_k$ included versus not included. }
    \label{fig:signvsnosign}
\end{figure}


In the context described above, we found that the model learned best when given the conductor $N_k$ as a feature. In this case, as well as in the cases where the model was given the factors of the conductor $\text{factor}(N_k)$ and the factors of the discriminant $\text{factor}(\Delta_k)$, the model reached 100\% accuracy. The fastest regime for reaching this level was with $N_k$, closely followed by $\text{factor}(\Delta_k)$ and closely followed  by  $\text{factor}(N_k)$. These findings are recorded in Figure~\ref{fig:three_extras}. Given the estimate \eqref{orderofnk} we are not surprised by this result, as $\frac{N_k}{\log (k)}$ essentially codifies the size of $n_k$, and $\log(k)$ is a relatively small quantity.
Moreover, we speculate that the model recovers $N_k$ from $\text{factor}(\Delta_k)$ and $\text{factor}(N_k)$ and
is unable to recover $N_k$ from $\Delta_k$ (due to the fact that factorization is difficult). Nevertheless, we find it puzzling that $\text{factor}(\Delta_k)$ leads to faster learning than $\text{factor}(N_k)$.
The other considered features, namely the rank and the $a_p$'s do not seem to influence the model's learning.
\begin{figure}
    \centering
    \includegraphics[scale = 0.5]{cmp_PB_3_6_67.pdf}
    \caption{Learning $r_k$ from $k$ and $\omega_k$ in three variants: Giving $N_k$, $\text{factor}(N_k)$ or $\text{factor}(\Delta_k)$.}
    \label{fig:three_extras}
\end{figure}


Finally, in Figure~\ref{fig:SnapshotofL2} we see that the values of $L(E_k,2)$ induced by the model's predictions have a marginal
distribution resembling the empirical distribution of the true values from Figure~\ref{fig:distrL_E_2}. This does not imply accurate pointwise prediction of $L(E_k,2)$.
 Interestingly,  the experiment given by $k, \omega_k$ and $\text{factor}(\Delta_k)$ seems to reproduce this distribution the best.

\begin{figure}
    \centering
    \includegraphics[scale = 0.5]{L2_snapshot_ep260.pdf}
    \caption{Distribution of the model guesses for $L(E_k,2)$ against the true values of $L(E_k,2)$, for $k$ in the validation set. This is for the experiment where  $k, \omega_k$ and $\text{factor}(\Delta_k)$ are given as input.}
    \label{fig:SnapshotofL2}
\end{figure}

\subsection{$p$-adic valuations of $n_k$ for $p \geq 5$.} These experiments were performed for $p=5,7$. They had $v_p(n_k)$ as output  with $k$ as a lone feature, and also with additional feature $\text{factor}(N_k)$, as well as additional feature $\text{factor}(\Delta_k)$ (only for $p=5$), all written in base $p$. In all cases
the model eventually predicts that none of the $n_k$ are multiple of $p$, thus leading to a greedy accuracy of
$\frac{p-2}{p-1}$. Therefore, we cannot extract any conclusions from these experiments.


\subsection{$3$-adic valuations of $n_k$}


For $p = 3$, we found that the model learned best when given $k$ and the factors of the discriminant $\mathrm{factor}(\Delta_k)$%
as features and $v_3(n_k)$ as output, all written in base $3$. We recorded the validation accuracy in Figure~\ref{fig:v3accuracy}. \begin{figure}
    \centering
    \includegraphics[scale = 0.5]{greedy_accuracy_p3_vs_epoch.pdf}
    \caption{The validation accuracy (i.e. the percentage of correct guesses by the model) as a function of the epochs.}
    \label{fig:v3accuracy}
\end{figure} From this figure, we learn that the accuracy maximizes near $51 \%$ around epoch $264$. Moreover, we see that the model starts overfitting from roughly epoch $300$ on. In this pre-overfitting phase, the model seems to almost always guess that $v_3(n_k) \in \{0, 1, 2, 3\}$ (this happens $99.99\%$ of the time, on average). The accuracy for $n_k$ with high $3$-valuations, $v_3(n_k) \in \{ 2, 3\}$ is quite low, see Figure~\ref{fig:acc_v3_three}. In this figure, we learn that validation accuracy of the model conditioning $v_3(n_k) = 1$ is nearly $100 \%$ all the time. Though this is a bit misleading:  from the earliest epochs the model simply guesses $v_3(n_k)=1$
for nearly every input, and so trivially classifies the true $v_3(n_k) = 1$ cases correctly (this also explains the starting plateau in Figure~\ref{fig:v3accuracy}). The surprising fact is that this near $100\%$ accuracy on $v_3 = 1$ persists even as the model learns to recognize $v_3 = 0$ cases, and the corresponding curve climbs.


\begin{figure}
    \centering
    \includegraphics[scale = 0.5]{exp19_acc_v3_three__2_.pdf}
    \caption{The validation accuracy, conditioned on $v_3(n_k)$}
    \label{fig:acc_v3_three}
\end{figure}

\subsubsection{Learning on different congruence classes}

As observed in the previous section, we have that $v_3(n_k) \geq 1$, whenever $k \equiv 0, \pm 5 \pmod 27$. We have that the model quite quickly learns this fact: $\mathbb{P}(\text{model predicts }v_3(n_k) = 0 \mid k \equiv 0, \pm 5 \pmod 27 )$ exceeds $1 \%$ only at $7$ epochs in the entire run. In fact, something more concerning happens: the model predicts almost always (i.e. $99.3 \%$ of the time in epoch range $100-300$) that $v_3(n_k) = 1$ for those $k$.  So, even though these congruence classes look more deterministic, their accuracy is in fact the lowest among all the congruence classes (their accuracy at epoch $264$ is just $41.5 \%$, compared to the global accuracy at this epoch being $50.9 \%$).

\subsubsection{Learning on the sets from Table \ref{tab:sets}}

Next we study how well the model learns the laws \eqref{eq:law-general} for $k \in A$. In Figure~\ref{fig:v3_three_A} we measure the validation accuracy assuming $v_3(n_k) = 0, 1$ or $>1$. \begin{figure}
    \centering
    \includegraphics[scale = 0.5]{acc_v3_three_A.pdf}
    \caption{The validation accuracy on $k \in A$, conditioned on $v_3(n_k)$.}
    \label{fig:v3_three_A}
\end{figure} From this figure, we learn that the accuracy for $v_3(n_k) > 1$ is still very poor, but that the accuracy assuming $v_3(n_k) = 0$ is $98.7 \%$. Thus, on the set $A$, the model is highly successful  in distinguishing $v_3(n_k) = 0$ from $v_3(n_k) = 1$. Figure~\ref{fig:acc-by-set} shows the model's validation accuracy, restricted to $v_3(n_k)\in\{0,1\}$, on each of the sets in Table~\ref{tab:sets}, both at the peak epoch and as a function of training. \begin{figure}[t]
\centering
\begin{minipage}[c]{0.30\textwidth}
  \centering\small
  \renewcommand{\arraystretch}{1.15}
  \begin{tabular}{@{}l r r@{}}
  \toprule
  set & $\#$ & acc.\\
  \midrule
  $A$   & $306$  & $99.3\%$\\
  $A'$  & $340$  & $97.4\%$\\
  $A''$ & $566$  & $95.1\%$\\
  \midrule
  $B$   & $1{,}133$ & $89.8\%$\\
  $B'$  & $977$  & $95.2\%$\\
  $B''$ & $937$  & $95.1\%$\\
  \midrule
  $C$   & $1{,}251$ & $88.6\%$\\
  $C'$  & $1{,}081$ & $93.8\%$\\
  $C''$ & $1{,}035$ & $93.6\%$\\
  \bottomrule
  \end{tabular}
\end{minipage}\hfill
\begin{minipage}[c]{0.66\textwidth}
  \centering
  \includegraphics[width=\linewidth]{acc_by_set_epochs.pdf}
\end{minipage}
\caption{Accuracy on $v_3(n_k)\in\{0,1\}$ over the sets of
Table~\ref{tab:sets}. Left: accuracy at the peak epoch~$264$, with $n$ the
number of validation $k$ in each set. Right: the same accuracies as a
function of the epochs.}
\label{fig:acc-by-set}
\end{figure} In this figure, we see that this conditional accuracy is the highest on the set $A$.  Although the parity laws remain
exact on some of the larger sets, these sets admit more varied factorizations
of $k(k^2-16)$. We speculate that this makes the relevant parity information
more difficult for the model to extract, leading to lower accuracy.
Removing the values $k \equiv 4 \pmod 8$, however --- passing from $B$ to $B'$ and from $C$ to $C'$ --- does increase this accuracy again.

Finally, we consider the complement of $C'' \cup \{ k \equiv 0, \pm 5 \pmod 27\}$, i.e. the values not covered by our largest set or the congruence classes above. As Figure~\ref{fig:outside} shows, the model accuracy is still increasing on the set. This leads us to suspect the model is still learning some arithmetic features, beyond the ones we already discussed.


\begin{figure}
    \centering
    \includegraphics[scale=0.5]{acc_outside_Cpp.pdf}
    \caption{The accuracy assuming $v_3(n_k) \in \{0,1\}$ outside $C''$ and $k \equiv 0, \pm 5 \pmod 27$.}
    \label{fig:outside}
\end{figure}


 After the experiment giving $\text{factor}(\Delta_k)$, the model learns the most if given the factors of the conductor
$\text{factor}(N_k)$. The
model learns faster from
$\text{factor}(N_k)$ than from $\text{factor}(\Delta_k)$. Giving both
$\text{factor}(N_k)$ and $\text{factor}(\Delta_k)$ together as features
 does not lead to a significant difference in learning compared to
 only giving $\text{factor}(N_k)$, but the learning is slower.
The model does not exhibit learning if given $N_k$ or $\Delta_k$ as features. We speculate that this is due to the difficulty in factorization.

\subsection{$2$-adic valuations of $n_k$}
For $p = 2$, we found that the model learned best when given $k$ and the factorization of the discriminant $\mathrm{factor}(\Delta_k)$ as input and $v_2(n_k)$ as output, all written in base $2$. The validation accuracy is recorded in Figure~\ref{fig:v2_val_acc}. We see that the model's overfitting starts happening (i.e. the validation accuracy decreases) around epochs $220 - 240$. From the figure we learn that the model climbs almost immediately to an accuracy of about $23 \% - 25 \%$ in the first few epochs and has a \emph{first} plateau there from roughly epoch $5$ to $75$. It then breaks through and starts climbing through the epochs $80-150$. After that it settles at a \emph{second} plateau with accuracy of roughly $36 \% - 38 \%$, maximizing at epoch $210$ with an accuracy $39.6 \%$. We analyze the behavior of the model at these plateaus.
\begin{figure}
    \centering
\includegraphics[scale = 0.5]{validation_accuracy.pdf}
    \caption{The validation accuracy (i.e. the percentage of correct guesses by the model) as a function of the epochs.}
    \label{fig:v2_val_acc}
\end{figure}


\subsubsection{Learning around the first plateau of Figure~\ref{fig:v2_val_acc}}
We first investigate whether the model learns our conjectural lower bound $\hat{v}(k) \leq v_2(n_k)$ from Conjecture~\ref{conj:low_bound}. As Figure~ \ref{fig:v2_lowbo} shows, from epoch 8 onward the model's predictions respect this bound for almost all $k$ in the validation set.
\begin{figure}
    \centering
    \includegraphics[scale=0.5]{lowerbound_proportion_vs_epoch__3_.pdf}
    \caption{The proportion of $k$ in the validation set for which the model predicts $\hat{v}(k) \leq v_2(n_k)$.}
    \label{fig:v2_lowbo}
\end{figure}

In particular, we compare our predictor $\hat{P}(k) = \mathrm{round}(\hat{v}(k) + s(k) + 1)$ (given in \eqref{eq:P}) with the model's prediction by computing the Pearson correlation coefficient between the two as a function of the epochs, see Figure~\ref{fig:pearson_pred}. We see that this correlation rises quickly and reaches a maximum of $r = 0.97$ around epochs $40-70$. As it leaves the first plateau, the correlation gradually drops to around $r \approx 0.88$ in the last epoch. Since the accuracy rises after the first plateau, but the correlation is dropping, this means that the model learns something beyond our predictor $\hat{P}(k)$.
For contrast, the correlation between the true valuations $v_2(n_k)$ and our predictor over the full dataset of $
249,999$ values is only $r = 0.78$. This behavior is illustrated in Figure~ \ref{fig:predictor_vs_model}, where we plot the confusion matrix of the model's prediction against our deterministic predictor $\hat{P}(k)$. It can be seen that the mass of the confusion matrix lies close to the diagonal. As the number of epochs grows further, the distribution spreads out from the diagonal even before overfitting. Together with the simultaneous increase in accuracy, this suggests that the model learns additional structure beyond $\hat{P}(k)$.%
Figure~\ref{fig:truth_vs_model} shows the analogous confusion matrix of the true valuation $v_2(n_k)$ against the model's prediction. Since the correlation between $\hat{P}(k)$
and the model's prediction is very close to $1$ throughout the first plateau, we conclude that during this phase the model has essentially learned to approximate $\hat{P}(k)$.
\begin{figure}[p]
    \centering

    \begin{subfigure}{0.95\textwidth}
        \centering
        \includegraphics[width=\textwidth,page=1]{confusion_pred_vs_Phat_snapshots.pdf}
        \caption{Model prediction vs. predictor $\hat P(k) = \mathrm{round}(\hat{v}(k) + s(k) + 1)$.}
        \label{fig:predictor_vs_model}
    \end{subfigure}

    \vspace{1em}

        \begin{subfigure}{0.95\textwidth}
        \centering
        \includegraphics[width=\textwidth,page=1]{confusion_true_vs_pred_snapshots.pdf}
        \caption{Model prediction vs. true $v_2(n_k)$.}
        \label{fig:truth_vs_model}

    \end{subfigure}

    \caption{Confusion-matrix snapshots across training epochs, including the Pearson correlation.}
    \label{fig:confusion_snapshots}
\end{figure}
\begin{figure}
    \centering
    \includegraphics[scale =0.5]{pearson_pred_vs_Phat_vs_epoch}
    \caption{The Pearson correlation between the model predictions and our predictor $\hat{P}(k)$ for $k$ in the validation set as a function of the epochs.}
\label{fig:pearson_pred}
\end{figure}
As guessing $\hat{P}(k)$ is correct only $24.3 \%$ of the time, and the model's accuracy on the first plateau is likewise about $23\% - 25 \%$, it appears that at this stage the model has learned nothing beyond this predictor.


\subsubsection{Learning around the second plateau of Figure~\ref{fig:v2_val_acc}} \label{sec:secondplateau}

For the second plateau, it becomes much less clear \emph{what} the model actually learned. But there are several hints, for instance,  the main improvement from the first to the second plateau is to detect when the lower bound in Conjecture \ref{conj:low_bound} is sharp, i.e. when $v_2(n_k) = \hat{v}(k)$. To be more precise, the proportion \[q_g := \mathbb{P}(\text{model is correct}\mid G(k) = v_2(n_k) - \hat{v}(k)= g) \] for $g = 0$ moves from $23.4\%$ at epoch $40$ to $76.9 \%$ at epoch $210$, see Figure~\ref{fig:correct_given_G_vs_epoch}. \begin{figure}[h]
    \centering
    \includegraphics[scale = 0.5]{correct_given_G_vs_epoch.pdf}
    \caption{The proportion of times the model is correct, given that $v_2(n_k) = \hat{v}(k)$ as a function of the epochs.}
    \label{fig:correct_given_G_vs_epoch}
\end{figure}
Additionally for $g = 1$, $q_g$ moves from $61.7 \%$ at epoch $40$ to $74 \%$ at epoch $210$. Here the high proportion at the first plateau can be explained from the fact the model learns to approximate $\hat{P}(k)$, which coincides with $\hat{v}(k) + 1$ (i.e. $g = 1$ here) quite often. For higher values of $g$, these improvements are much smaller, see Table \ref{tab:correct_by_G}.
\begin{table}[ht]
  \centering
  \begin{subtable}[t]{0.48\linewidth}
    \centering
    \begin{tabular}{ccc}
      \toprule
      $g$ & $\mathbb{P}(G(k)=g)$ & $\mathbb{P}(\text{correct}\mid G(k)=g)$ \\
      \midrule
      $0$ & $22.4\%$ & $23.4\%$ \\
      $1$ & $24.5\%$ & $61.7\%$ \\
      $2$ & $21.3\%$ & $21.5\%$ \\
      $3$ & $13.5\%$ & $0.2\%$  \\
      $4$ & $8.4\%$  & $0.0\%$  \\
      \bottomrule
    \end{tabular}
    \caption{Epoch $40$ (first plateau).}
    \label{tab:correct_epoch40}
  \end{subtable}
  \hfill
  \begin{subtable}[t]{0.48\linewidth}
    \centering
    \begin{tabular}{ccc}
      \toprule
      $g$ & $\mathbb{P}(G(k)=g)$ & $\mathbb{P}(\text{correct}\mid G(k)=g)$ \\
      \midrule
      $0$ & $22.4\%$ & $76.9\%$ \\
      $1$ & $24.5\%$ & $74.0\%$ \\
      $2$ & $21.3\%$ & $17.9\%$ \\
      $3$ & $13.5\%$ & $3.5\%$  \\
      $4$ & $8.4\%$  & $0.4\%$  \\
      \bottomrule
    \end{tabular}
    \caption{Epoch $210$ (second plateau).}
    \label{tab:correct_epoch210}
  \end{subtable}
  \caption{The proportion of times the model guesses correctly, under the assumption
    that the value $G(k) = v_2(n_k) - \hat{v}(k)$ is fixed, together with the
    distribution of $G(k)$. Comparing the first plateau (A) with the second plateau (B). This is for all $k$ in the validation set.}
  \label{tab:correct_by_G}
\end{table}


 The optimal learning with   $\text{factor}(\Delta_k)$ as a feature should be interpreted in light of equations  \eqref{eq:hatvdef} and \eqref{eq:v2approx}, as the factors of the discriminant encode not only the primes of bad reduction, but they also make a distinction between those dividing $k$ and those dividing $k^2-16$.
Among the other experiments that were considered, providing the  factors of the conductor
$\text{factor}(N_k)$ as a feature leads to substantial learning, but less than giving $\text{factor}(\Delta_k)$. We interpret this phenomenon as providing the primes of bad reduction, but not making the distinction among those that divide $k$ from those that divide $k^2-16$.  Giving both
$\text{factor}(N_k)$ and $\text{factor}(\Delta_k)$ together as features leads to slightly worse learning than giving
$\text{factor}(\Delta_k)$ alone, but significantly better than giving $\text{factor}(N_k)$ alone. We speculate that the model has some trouble when provided with too much information. We found no substantial learning when introducing $c_2$, $N_k$, or $\Delta_k$ as features. For the last two cases, this is consistent with the difficulty in factorization.


\section{Conclusions and future directions}
\label{sec:conclusions}

In this work, we have investigated the rational factors $r_k$ appearing
conjecturally in Boyd's identities \eqref{eq:conjecture}
by combining a large-scale numerical study with machine-learning experiments.
Rather than revealing a single formula for $r_k$, our explorations reveal
several layers of arithmetic structure. The sign of $r_k$ is governed by the
root number, while the conductor determines the order of magnitude of
$n_k=1/r_k$.
At the same time, the divisibility properties of $n_k$ are strongly influenced
by the local arithmetic of $E_k$.

One of the most striking conclusions of the statistical analysis is the sharp
difference between the behavior at different primes. For primes $p\geq 5$, the
positive valuations appear to follow the remarkably simple law $\mathbb{P}\bigl(v_p(n_k)=r\bigr)=p^{-r}$ for $r\geq 1,$
although the deviation observed at $p=5$ suggests that this model may require
some refinement. The primes $2$ and $3$ exhibit considerably richer
arithmetic. At $3$, congruence conditions on $k$ interact with the parity of
the number of bad primes congruent to $1$ modulo $3$, producing laws that are
essentially deterministic on large subfamilies. At $2$, the valuation is well
approximated by a weighted count of the odd primes dividing $k$ and
$k^2-16$, together with corrections arising from the reduction at $2$ and
from the $2$-adic proximity of $k$ to $\pm8$. After these deterministic contributions are removed, the remaining term is
closely modeled by a fixed negative-binomial distribution. Although this last
observation remains heuristic, the lower bound of Conjecture~\ref{conj:low_bound}
provides a concrete arithmetic statement suggested by the data.

The machine-learning experiments clarify both the possibilities and the
limitations of the present approach. Predicting $r_k$ exactly proved to be too
difficult, but the model learned its size very effectively when the conductor, or factorizations from which it can be recovered, were included among the inputs. The experiments
therefore confirm that the choice and representation of arithmetic features
are crucial: the factorizations of the conductor and discriminant are much
more informative for the valuation problems than the corresponding unfactored
integers. By contrast, the experiments for $p=5$ and $p=7$ collapsed to the
most frequent output and did not reveal additional structure, illustrating
the difficulty caused by highly unbalanced valuation distributions.

The results at $2$ and $3$ are more encouraging. For the $3$-adic valuation,
the model recognizes several of the congruence and parity phenomena found in
the statistical analysis and continues to improve even on values of $k$ not
covered by the largest families considered in Section~3. More surprisingly,
in the $2$-adic experiment the evolution of the model appears to have a
mathematical interpretation. During the first plateau, its predictions are
very close to the deterministic predictor $\widehat P(k)$ derived from our
statistical model. During the second plateau, however, its accuracy increases
while its correlation with $\widehat P(k)$ decreases. In particular, the
model becomes substantially better at recognizing when the lower bound
$\widehat v(k)\leq v_2(n_k)$ is sharp. This suggests that the network detects
additional arithmetic information that is not yet incorporated into our
description of $v_2(n_k)$.

These observations lead to several natural questions.
 The most important
is to find a conceptual explanation for the arithmetic of $n_k$. This includes
explaining its near-integrality, proving or refining the proposed distribution
for $v_p(n_k)$ when $p\geq5$, and understanding the congruence and parity laws
at $3$. At the prime $2$, one would like to prove
Conjecture~\ref{conj:low_bound}, characterize the values of $k$ for which
equality holds, and explain the apparent negative-binomial distribution of
the excess. The conditional formula discussed in
Remark~\ref{rem:general-BK} suggests a possible framework for such
an explanation.  It would be interesting to develop arithmetic
statistics for the families of Tate-twisted Bloch--Kato Selmer
groups and regulator indices occurring in that formula.  Existing
probabilistic models for ordinary Selmer groups, such as those
surveyed by Poonen \cite{Poonen}, may provide useful guidance, although the groups
appearing here are different and the regulator-index term must also
be modeled.

For $p\geq5$, the proposed law
\[
\mathbb P\bigl(v_p(n_k)=r\bigr)=p^{-r}
\]
says that, after the first factor of $p$ occurs, the higher
valuation follows the usual geometric law.  The main question
is therefore the increased probability
\[
\mathbb P(v_p(n_k) \geq 1)=\frac{1}{p-1}.
\]
It would be interesting to determine whether this
enhancement arises from the local Tamagawa factors, the
Bloch--Kato Selmer group, the regulator index, or an interaction
among these terms.  Explicit approaches to $p$-adic regulators,
such as those developed by Coleman--de Shalit~\cite{Coleman-de-Shalit}, Besser ~\cite{Besser}, and
Asakura--Chida ~\cite{AsakuraChida}, may be useful in investigating the regulator
contribution. Ultimately, these phenomena should be related to the regulator
interpretations underlying Boyd's conjectures and to the local arithmetic of
the elliptic curves $E_k$.


The machine-learning results also suggest  directions for future
experiments.
Instead of predicting $v_2(n_k)$ directly, one could train a model
to predict the excess $G(k)=v_2(n_k)-\widehat v(k),$
or simply to classify the cases in which $G(k)=0$. We attempted to use the model to identify the cases in which $G(k)=0$, but did not obtain a useful classification. Interpreting such a
classifier, testing additional arithmetic features, and searching for
subfamilies on which it performs unusually well may lead to a refinement of
the present conjectures. For the primes $p\geq5$, the model tends to collapse to the trivial prediction $v_p(n_k) = 0$, since this is the most common class and the loss function only distinguishes correct from incorrect predictions. A class-balanced loss, which assigns greater weight to the rarer positive valuations, or a probabilistic model predicting the full distribution of $v_p(n_k)$, may help prevent this collapse.%
A first natural extension, taking $k=\sqrt m$ for $m\in\mathbb Z$ of either sign, is
carried out in Section~\ref{sec:sqrtm}: $r_{\sqrt m}$ is again rational with
near-integral reciprocal and the valuation statistics for $p\neq3$ extend to this setting, but the finer $3$-adic parity law of
Section~\ref{ssec:partitylaw} breaks down. Finding its correct generalization to
nonsquare $m$ would be very interesting, as would
testing whether these patterns persist in other one-parameter families of Mahler
measures associated with elliptic curves.

Thus, although the experiments do not provide a formula for $r_k$, they show
that these rational factors are far from arithmetically unstructured. The
combination of large-scale computation, statistical analysis, and machine
learning has produced concrete conjectures and has identified arithmetic phenomena that are not evident
from the currently known identities alone. The next step is
to turn these empirical regularities into arithmetic explanations and,
eventually, proofs.


\clearpage
\section{Appendix: Tables and Figures}
\label{sec:appendix}


\begin{table}[htbp]
\centering
\scriptsize
\begin{tabular}{@{}c|cccccccccc||c@{}}
$(p,r)$ & $[0,25K]$ & $[0,50K]$ & $[0,75K]$ & $[0,100K]$ & $[0,125K]$ & $[0,150K]$ & $[0,175K]$ & $[0,200K]$ & $[0,225K]$ & $[0,250K]$ & $p^{-r}$ \\
\hline
$(5,1)$ & 0.209 & 0.208 & 0.205 & 0.205 & 0.204 & 0.203 & 0.202 & 0.202 & 0.202 & 0.202 & 0.200 \\
$(5,2)$ & 0.044 & 0.044 & 0.044 & 0.044 & 0.043 & 0.043 & 0.043 & 0.043 & 0.043 & 0.043 & 0.040 \\
$(5,3)$ & 0.009 & 0.009 & 0.009 & 0.009 & 0.009 & 0.009 & 0.008 & 0.008 & 0.009 & 0.009 & 0.008 \\
$(5,4)$ & 0.002 & 0.002 & 0.002 & 0.002 & 0.002 & 0.002 & 0.002 & 0.002 & 0.002 & 0.002 & 0.002 \\
\hline\hline
$(7,1)$ & 0.142 & 0.143 & 0.144 & 0.144 & 0.143 & 0.142 & 0.142 & 0.142 & 0.142 & 0.142 & 0.143 \\
$(7,2)$ & 0.021 & 0.021 & 0.020 & 0.020 & 0.020 & 0.020 & 0.021 & 0.020 & 0.021 & 0.021 & 0.020 \\
$(7,3)$ & 0.003 & 0.003 & 0.003 & 0.003 & 0.003 & 0.003 & 0.003 & 0.003 & 0.003 & 0.003 & 0.003 \\
\hline\hline
$(11,1)$ & 0.088 & 0.090 & 0.091 & 0.090 & 0.091 & 0.091 & 0.091 & 0.090 & 0.091 & 0.091 & 0.091 \\
$(11,2)$ & 0.008 & 0.009 & 0.008 & 0.008 & 0.008 & 0.008 & 0.008 & 0.008 & 0.009 & 0.009 & 0.008 \\
$(11,3)$ & 0.001 & 0.001 & 0.001 & 0.001 & 0.001 & 0.001 & 0.001 & 0.001 & 0.001 & 0.001 & 0.001 \\
\hline\hline
$(13,1)$ & 0.074 & 0.075 & 0.076 & 0.076 & 0.076 & 0.076 & 0.076 & 0.075 & 0.076 & 0.076 & 0.077 \\
$(13,2)$ & 0.006 & 0.006 & 0.006 & 0.006 & 0.006 & 0.006 & 0.006 & 0.006 & 0.006 & 0.006 & 0.006 \\
$(13,3)$ & 0.000 & 0.001 & 0.001 & 0.000 & 0.000 & 0.000 & 0.000 & 0.000 & 0.000 & 0.000 & 0.000 \\
\hline\hline
$(17,1)$ & 0.059 & 0.058 & 0.059 & 0.059 & 0.059 & 0.059 & 0.059 & 0.059 & 0.059 & 0.059 & 0.059 \\
$(17,2)$ & 0.003 & 0.003 & 0.003 & 0.003 & 0.003 & 0.003 & 0.003 & 0.003 & 0.003 & 0.003 & 0.003 \\
\hline\hline
$(19,1)$ & 0.051 & 0.052 & 0.051 & 0.051 & 0.052 & 0.051 & 0.052 & 0.052 & 0.052 & 0.052 & 0.053 \\
$(19,2)$ & 0.003 & 0.003 & 0.003 & 0.003 & 0.003 & 0.003 & 0.003 & 0.003 & 0.003 & 0.003 & 0.003 \\
\hline\hline
$(23,1)$ & 0.043 & 0.043 & 0.043 & 0.043 & 0.043 & 0.043 & 0.043 & 0.043 & 0.043 & 0.043 & 0.043 \\
$(23,2)$ & 0.001 & 0.002 & 0.002 & 0.002 & 0.002 & 0.002 & 0.002 & 0.002 & 0.002 & 0.002 & 0.002 \\
\hline\hline
$(29,1)$ & 0.036 & 0.035 & 0.035 & 0.035 & 0.035 & 0.035 & 0.035 & 0.035 & 0.035 & 0.035 & 0.034 \\
\hline\hline
$(31,1)$ & 0.033 & 0.033 & 0.033 & 0.033 & 0.033 & 0.033 & 0.033 & 0.033 & 0.033 & 0.033 & 0.032 \\
\hline\hline
$(37,1)$ & 0.025 & 0.026 & 0.026 & 0.026 & 0.027 & 0.027 & 0.027 & 0.027 & 0.027 & 0.027 & 0.027 \\
\end{tabular}
\caption{Observed cumulative proportions of $k$ satisfying $v_p(n_k)=r$ for the pairs $(p,r)$ listed in the first column, across closed intervals of the form $[0,x]$, where $K$ denotes $1000$, together with the reference value $p^{-r}$.}
\label{tab:v5_v7_v11_v13_v17_v19_v23_v29_v31_v37_statistics_cumulative}
\end{table}


\begin{table}
\centering
\scriptsize
\begin{tabular}{@{}c|cccccccccc||c@{}}
$(p,r)$
& $[0,25K]$ & $[0,50K]$ & $[0,75K]$ & $[0,100K]$ & $[0,125K]$
& $[0,150K]$ & $[0,175K]$ & $[0,200K]$ & $[0,225K]$ & $[0,250K]$
& $p^{-r}$ \\
\hline
$(3,1)$ & 0.325 & 0.322 & 0.323 & 0.322 & 0.322 & 0.321 & 0.322 & 0.322 & 0.321 & 0.322 & 0.333 \\
$(3,2)$ & 0.231 & 0.233 & 0.234 & 0.234 & 0.234 & 0.234 & 0.234 & 0.234 & 0.235 & 0.234 & 0.111 \\
$(3,3)$ & 0.103 & 0.106 & 0.107 & 0.108 & 0.108 & 0.110 & 0.110 & 0.110 & 0.110 & 0.110 & 0.037 \\
$(3,4)$ & 0.037 & 0.041 & 0.042 & 0.042 & 0.042 & 0.042 & 0.042 & 0.042 & 0.043 & 0.043 & 0.012 \\
$(3,5)$ & 0.014 & 0.014 & 0.014 & 0.014 & 0.015 & 0.015 & 0.015 & 0.015 & 0.015 & 0.015 & 0.004 \\
$(3,6)$ & 0.005 & 0.005 & 0.005 & 0.005 & 0.005 & 0.005 & 0.005 & 0.005 & 0.005 & 0.005 & 0.001 \\
\end{tabular}
\caption{Observed cumulative proportions of $k$ satisfying $v_3(n_k)=r$ for the pairs $(3,r)$ listed in the first column, across closed intervals of the form $[0,x]$, where $K$ denotes $1000$.}
\label{tab:v3_statistics_cumulative}
\end{table}


\begin{table}[htbp]
\centering
\scriptsize
\begin{tabular}{@{}cc|cccccccccc@{}}
$(p,r)$ & $k\pmod 3$ & $[0,25K]$ & $[0,50K]$ & $[0,75K]$ & $[0,100K]$ & $[0,125K]$ & $[0,150K]$ & $[0,175K]$ & $[0,200K]$ & $[0,225K]$ & $[0,250K]$ \\
\hline
$(3,1)$ & all & 0.718 & 0.724 & 0.728 & 0.728 & 0.728 & 0.730 & 0.731 & 0.731 & 0.732 & 0.732 \\
 & $0$ & 0.754 & 0.755 & 0.757 & 0.757 & 0.757 & 0.758 & 0.758 & 0.758 & 0.758 & 0.758 \\
 & $1$ & 0.698 & 0.708 & 0.714 & 0.713 & 0.713 & 0.714 & 0.715 & 0.717 & 0.718 & 0.719 \\
 & $2$ & 0.700 & 0.708 & 0.713 & 0.715 & 0.716 & 0.717 & 0.718 & 0.719 & 0.719 & 0.720 \\
\hline
$(3,2)$ & all & 0.392 & 0.401 & 0.405 & 0.406 & 0.407 & 0.408 & 0.409 & 0.410 & 0.411 & 0.411 \\
 & $0$ & 0.432 & 0.436 & 0.440 & 0.442 & 0.442 & 0.444 & 0.444 & 0.445 & 0.444 & 0.445 \\
 & $1$ & 0.375 & 0.385 & 0.389 & 0.389 & 0.389 & 0.390 & 0.391 & 0.393 & 0.394 & 0.394 \\
 & $2$ & 0.370 & 0.382 & 0.384 & 0.388 & 0.389 & 0.391 & 0.391 & 0.392 & 0.394 & 0.393 \\
\hline
$(3,3)$ & all & 0.161 & 0.169 & 0.171 & 0.172 & 0.173 & 0.175 & 0.175 & 0.176 & 0.176 & 0.176 \\
 & $0$ & 0.187 & 0.193 & 0.194 & 0.196 & 0.195 & 0.196 & 0.197 & 0.197 & 0.198 & 0.198 \\
 & $1$ & 0.149 & 0.156 & 0.161 & 0.163 & 0.163 & 0.165 & 0.166 & 0.166 & 0.166 & 0.166 \\
 & $2$ & 0.148 & 0.156 & 0.157 & 0.159 & 0.161 & 0.163 & 0.163 & 0.164 & 0.165 & 0.165 \\
\hline
\end{tabular}
\caption{Observed cumulative proportions of $k$ satisfying $v_3(n_k)\ge r$ for the pairs $(3,r)$ listed in the first column, split according to congruence classes of $k \pmod 3$, across closed intervals of the form $[0,x]$, where $K$ denotes $1000$. }
\label{tab:v3_geq_mod3_statistics_congruence_cumulative}
\end{table}


\begin{figure}[htbp]
    \centering
    \includegraphics[scale=0.7]{histogram_all_v3_0_250K.pdf}

    \caption{Histogram of the $3$-adic valuations $v_3(n_k)$ for $1 \leq k \leq 250K$ such that  $v_3(n_k) \geq 0$, where $K=1000$.}
    \label{fig:histogram_v3_0_250K}
\end{figure}

\begin{table}[htbp]
\centering
\scriptsize
\begin{tabular}{@{}c|cccccccccc@{}}
$(p,r)$ & $[0,25K]$ & $[0,50K]$ & $[0,75K]$ & $[0,100K]$ & $[0,125K]$ & $[0,150K]$ & $[0,175K]$ & $[0,200K]$ & $[0,225K]$ & $[0,250K]$ \\
\hline
$(2,1)$ & 0.000 & 0.000 & 0.000 & 0.000 & 0.000 & 0.000 & 0.000 & 0.000 & 0.000 & 0.000 \\
$(2,2)$ & 0.002 & 0.002 & 0.001 & 0.001 & 0.001 & 0.001 & 0.001 & 0.001 & 0.001 & 0.001 \\
$(2,3)$ & 0.011 & 0.008 & 0.007 & 0.006 & 0.005 & 0.005 & 0.005 & 0.005 & 0.004 & 0.004 \\
$(2,4)$ & 0.028 & 0.021 & 0.018 & 0.016 & 0.015 & 0.014 & 0.014 & 0.013 & 0.012 & 0.012 \\
$(2,5)$ & 0.053 & 0.043 & 0.038 & 0.035 & 0.033 & 0.031 & 0.030 & 0.029 & 0.028 & 0.027 \\
$(2,6)$ & 0.089 & 0.075 & 0.068 & 0.064 & 0.061 & 0.058 & 0.056 & 0.055 & 0.053 & 0.052 \\
$(2,7)$ & 0.131 & 0.115 & 0.107 & 0.102 & 0.097 & 0.093 & 0.091 & 0.089 & 0.087 & 0.085 \\
$(2,8)$ & 0.147 & 0.140 & 0.133 & 0.129 & 0.126 & 0.123 & 0.121 & 0.119 & 0.117 & 0.115 \\
$(2,9)$ & 0.148 & 0.147 & 0.146 & 0.144 & 0.143 & 0.141 & 0.140 & 0.138 & 0.137 & 0.136 \\
$(2,10)$ & 0.125 & 0.134 & 0.137 & 0.137 & 0.137 & 0.138 & 0.138 & 0.138 & 0.138 & 0.138 \\
$(2,11)$ & 0.095 & 0.108 & 0.114 & 0.118 & 0.121 & 0.123 & 0.124 & 0.125 & 0.125 & 0.126 \\
$(2,12)$ & 0.067 & 0.080 & 0.086 & 0.090 & 0.093 & 0.096 & 0.098 & 0.100 & 0.102 & 0.103 \\
$(2,13)$ & 0.045 & 0.053 & 0.058 & 0.062 & 0.065 & 0.068 & 0.070 & 0.072 & 0.074 & 0.075 \\
$(2,14)$ & 0.026 & 0.033 & 0.037 & 0.040 & 0.042 & 0.044 & 0.045 & 0.047 & 0.048 & 0.050 \\
$(2,15)$ & 0.015 & 0.019 & 0.021 & 0.024 & 0.026 & 0.028 & 0.029 & 0.031 & 0.031 & 0.032 \\
$(2,16)$ & 0.007 & 0.010 & 0.013 & 0.014 & 0.015 & 0.016 & 0.017 & 0.018 & 0.018 & 0.019 \\
$(2,17)$ & 0.004 & 0.006 & 0.007 & 0.008 & 0.009 & 0.009 & 0.010 & 0.010 & 0.011 & 0.011 \\
$(2,18)$ & 0.003 & 0.003 & 0.004 & 0.005 & 0.005 & 0.005 & 0.005 & 0.006 & 0.006 & 0.006 \\
$(2,19)$ & 0.002 & 0.002 & 0.002 & 0.002 & 0.003 & 0.003 & 0.003 & 0.003 & 0.003 & 0.003 \\
$(2,20)$ & 0.001 & 0.001 & 0.001 & 0.001 & 0.001 & 0.001 & 0.002 & 0.002 & 0.002 & 0.002 \\
\end{tabular}
\caption{Observed cumulative proportions of $k$ satisfying  $v_2(n_k)=r$ for the pairs $(2,r)$ listed in the first column, across closed intervals of the form $[0,x]$, where $K$ denotes $1000$.}
\label{tab:v2_statistics_cumulative}
\end{table}


\begin{table}[htbp]
\centering
\scriptsize
\begin{tabular}{@{}cc|cccccccccc@{}}
$(p,r)$ & $k\pmod 7$ & $[0,25K]$ & $[0,50K]$ & $[0,75K]$ & $[0,100K]$ & $[0,125K]$ & $[0,150K]$ & $[0,175K]$ & $[0,200K]$ & $[0,225K]$ & $[0,250K]$ \\
\hline
$(2,6)$ & all & 0.905 & 0.926 & 0.935 & 0.941 & 0.945 & 0.948 & 0.951 & 0.953 & 0.954 & 0.956 \\
 & $0$ & 0.921 & 0.938 & 0.947 & 0.952 & 0.956 & 0.958 & 0.961 & 0.963 & 0.964 & 0.965 \\
 & $1$ & 0.863 & 0.892 & 0.905 & 0.914 & 0.919 & 0.923 & 0.927 & 0.930 & 0.933 & 0.935 \\
 & $2$ & 0.870 & 0.896 & 0.908 & 0.914 & 0.919 & 0.923 & 0.927 & 0.930 & 0.932 & 0.934 \\
 & $3$ & 0.971 & 0.982 & 0.987 & 0.989 & 0.990 & 0.991 & 0.992 & 0.992 & 0.993 & 0.993 \\
 & $4$ & 0.972 & 0.981 & 0.985 & 0.988 & 0.989 & 0.990 & 0.991 & 0.991 & 0.992 & 0.992 \\
 & $5$ & 0.865 & 0.896 & 0.909 & 0.916 & 0.921 & 0.926 & 0.928 & 0.931 & 0.933 & 0.935 \\
 & $6$ & 0.872 & 0.896 & 0.908 & 0.916 & 0.922 & 0.925 & 0.928 & 0.931 & 0.933 & 0.935 \\
\hline
$(2,7)$ & all & 0.816 & 0.851 & 0.867 & 0.877 & 0.884 & 0.890 & 0.894 & 0.898 & 0.901 & 0.903 \\
 & $0$ & 0.835 & 0.870 & 0.888 & 0.896 & 0.902 & 0.907 & 0.911 & 0.914 & 0.917 & 0.919 \\
 & $1$ & 0.754 & 0.800 & 0.818 & 0.833 & 0.841 & 0.848 & 0.854 & 0.859 & 0.862 & 0.865 \\
 & $2$ & 0.767 & 0.805 & 0.823 & 0.833 & 0.841 & 0.848 & 0.853 & 0.859 & 0.863 & 0.866 \\
 & $3$ & 0.917 & 0.942 & 0.952 & 0.958 & 0.962 & 0.964 & 0.967 & 0.969 & 0.970 & 0.971 \\
 & $4$ & 0.925 & 0.944 & 0.954 & 0.959 & 0.962 & 0.965 & 0.967 & 0.969 & 0.970 & 0.972 \\
 & $5$ & 0.754 & 0.798 & 0.818 & 0.832 & 0.842 & 0.849 & 0.853 & 0.858 & 0.862 & 0.865 \\
 & $6$ & 0.758 & 0.799 & 0.817 & 0.831 & 0.841 & 0.847 & 0.853 & 0.857 & 0.862 & 0.864 \\
\hline
\end{tabular}
\caption{Observed cumulative proportions of $k$ satisfying $v_2(n_k)\ge r$ for the pairs $(2,r)$ listed in the first column, split according to congruence classes of $k \pmod 7$, across closed intervals of the form $[0,x]$, where $K$ denotes $1000$.
\label{tab:v2_geq_mod7_statistics_congruence_cumulative}}
\end{table}


\begin{figure}[htbp]
    \centering
    \includegraphics[scale=0.7]{histogram_all_v2_0_250K.pdf}

    \caption{Histogram of the $2$-adic valuations $v_2(n_k)$ for $1 \leq k \leq 250K$ such that $v_2(n_k) \geq 0$, where $K=1000$.}
    \label{fig:histogram_v2_0_250K}
\end{figure}


\begin{table}[htbp]
\centering\small
\setlength{\tabcolsep}{4.5pt}
\begin{tabular}{r l l c c c c}
\toprule
$k$ & $k-4$ & $k+4$ & $W(k)$ & $\hat v(k)$ & $v_2(n_k)$ & $v_2(n_k)-\hat v(k)$\\
\midrule
$321559$  & $3\cdot5\cdot13\cdot17\cdot97$           & $11\cdot23\cdot31\cdot41$            & $21$ & $17$ & $19$ & $+2$\\
$400201$  & $3\cdot7\cdot17\cdot19\cdot59$           & $5\cdot13\cdot47\cdot131$           & $21$ & $17$ & $20$ & $+3$\\
$440891$  & $23\cdot29\cdot661$                      & $3\cdot5\cdot7\cdot13\cdot17\cdot19$ & $21$ & $17$ & $19$ & $+2$\\
$471709$  & $3\cdot5\cdot13\cdot41\cdot59$           & $11\cdot19\cdot37\cdot61$           & $21$ & $17$ & $21$ & $+4$\\
$474415$  & $3\cdot7\cdot19\cdot29\cdot41$           & $11\cdot17\cdot43\cdot59$           & $21$ & $17$ & $18$ & $+1$\\
$519589$  & $3\cdot5\cdot11\cdot47\cdot67$           & $19\cdot23\cdot29\cdot41$           & $21$ & $17$ & $20$ & $+3$\\
$572029$  & $3\cdot5^{2}\cdot29\cdot263$             & $7\cdot11\cdot17\cdot19\cdot23$     & $21$ & $17$ & $22$ & $+5$\\
$729491$  & $11\cdot17\cdot47\cdot83$                & $3^{2}\cdot5\cdot13\cdot29\cdot43$  & $21$ & $17$ & $21$ & $+4$\\
$889291$  & $3\cdot7\cdot17\cdot47\cdot53$           & $5\cdot11\cdot19\cdot23\cdot37$     & $23$ & $19$ & $21$ & $+2$\\
$1008011$ & $7\cdot11\cdot13\cdot19\cdot53$          & $3\cdot5\cdot17\cdot59\cdot67$      & $22$ & $18$ & $20$ & $+2$\\
$1394701$ & $3\cdot17\cdot23\cdot29\cdot41$          & $5\cdot13\cdot43\cdot499$           & $22$ & $18$ & $21$ & $+3$\\
$1601149$ & $3^{2}\cdot5\cdot7\cdot13\cdot17\cdot23$ & $83\cdot101\cdot191$                & $22$ & $18$ & $21$ & $+3$\\
$1847751$ & $11\cdot17\cdot41\cdot241$              & $5\cdot7\cdot13\cdot31\cdot131$         & $22$ & $18$ & $23$ & $+5$\\
$1944939$ & $5\cdot19\cdot59\cdot347$               & $7\cdot11\cdot13\cdot29\cdot67$         & $22$ & $18$ & $22$ & $+4$\\
\bottomrule
\end{tabular}
\caption{Prime-rich odd values of $k$ in $2.5\times10^{5}<k<2\times10^{6}$ chosen to
maximise $W(k):=\omega_{\mathrm{odd}}(k)+2\,\omega_{\mathrm{odd}}(k^2-16)$; the maximum
$W=23$ is attained uniquely at $k=889291$. For odd $k$ the curve $E_k$ is semistable at
$2$, so $c(k)=4$ and $\hat v(k)=W(k)-4$. The table includes every odd $k$ in the window
with $W\ge22$, together with a selection of those with $W=21$. In every case
$v_2(n_k)\ge\hat v(k)$. }
\label{tab:largeW-odd}
\end{table}


\begin{table}[htbp]
\centering\small
\setlength{\tabcolsep}{4.5pt}
\begin{tabular}{r l l c c c c}
\toprule
$k$ & $k-4$ & $k+4$ & $W(k)$ & $\hat v(k)$ & $v_2(n_k)$ & $v_2(n_k)-\hat v(k)$\\
\midrule
$940474$  & $2\cdot3\cdot5\cdot23\cdot29\cdot47$     & $2\cdot7\cdot11\cdot31\cdot197$         & $21$ & $19$ & $21$ & $+2$\\
$1069814$ & $2\cdot5\cdot7\cdot17\cdot29\cdot31$     & $2\cdot3\cdot37\cdot61\cdot79$          & $21$ & $19$ & $27$ & $+8$\\
$1195674$ & $2\cdot5\cdot7\cdot19\cdot29\cdot31$     & $2\cdot11\cdot17\cdot23\cdot139$        & $21$ & $19$ & $20$ & $+1$\\
$1204606$ & $2\cdot3\cdot7\cdot23\cdot29\cdot43$     & $2\cdot5\cdot11\cdot47\cdot233$         & $21$ & $19$ & $22$ & $+3$\\
$1234874$ & $2\cdot5\cdot7\cdot13\cdot23\cdot59$     & $2\cdot3\cdot29\cdot47\cdot151$         & $21$ & $19$ & $22$ & $+3$\\
$1258786$ & $2\cdot3\cdot7\cdot17\cdot41\cdot43$     & $2\cdot5\cdot13\cdot23\cdot421$         & $21$ & $19$ & $21$ & $+2$\\
$1286666$ & $2\cdot13\cdot17\cdot41\cdot71$          & $2\cdot3\cdot5\cdot7\cdot11\cdot557$    & $21$ & $19$ & $21$ & $+2$\\
$1540136$ & $2^{2}\cdot11\cdot17\cdot29\cdot71$      & $2^{2}\cdot3\cdot5\cdot7\cdot19\cdot193$ & $21$ & $19$ & $19$ & $+0$\\
$1587226$ & $2\cdot3^{3}\cdot7\cdot13\cdot17\cdot19$ & $2\cdot5\cdot23\cdot67\cdot103$         & $21$ & $19$ & $20$ & $+1$\\
$1614826$ & $2\cdot3\cdot11\cdot43\cdot569$          & $2\cdot5\cdot7\cdot17\cdot23\cdot59$    & $21$ & $19$ & $20$ & $+1$\\
$1653526$ & $2\cdot3\cdot13\cdot17\cdot29\cdot43$    & $2\cdot5\cdot37\cdot41\cdot109$         & $21$ & $19$ & $22$ & $+3$\\
$1669906$ & $2\cdot3\cdot13\cdot79\cdot271$          & $2\cdot5\cdot11\cdot17\cdot19\cdot47$   & $21$ & $19$ & $20$ & $+1$\\
$1700846$ & $2\cdot11\cdot13\cdot19\cdot313$         & $2\cdot3\cdot5^{2}\cdot17\cdot23\cdot29$ & $21$ & $19$ & $21$ & $+2$\\
$1703464$ & $2^{2}\cdot3\cdot5\cdot11\cdot29\cdot89$ & $2^{2}\cdot13\cdot17\cdot41\cdot47$     & $21$ & $19$ & $26$ & $+7$\\
$1717166$ & $2\cdot41\cdot43\cdot487$                & $2\cdot3\cdot5\cdot7\cdot13\cdot17\cdot37$ & $21$ & $19$ & $24$ & $+5$\\
$1722886$ & $2\cdot3\cdot7\cdot17\cdot19\cdot127$    & $2\cdot5\cdot13\cdot29\cdot457$         & $21$ & $19$ & $21$ & $+2$\\
$1762094$ & $2\cdot5\cdot11\cdot83\cdot193$          & $2\cdot3\cdot13\cdot19\cdot29\cdot41$   & $21$ & $19$ & $24$ & $+5$\\
$1774630$ & $2\cdot3\cdot7\cdot29\cdot31\cdot47$     & $2\cdot23\cdot173\cdot223$             & $21$ & $19$ & $22$ & $+3$\\
$1791566$ & $2\cdot17\cdot23\cdot29\cdot79$          & $2\cdot3\cdot5\cdot11\cdot61\cdot89$    & $21$ & $19$ & $24$ & $+5$\\
$1812794$ & $2\cdot5\cdot7\cdot19\cdot29\cdot47$     & $2\cdot3^{2}\cdot13\cdot61\cdot127$     & $21$ & $19$ & $19$ & $+0$\\
$1885594$ & $2\cdot3^{2}\cdot5\cdot7\cdot41\cdot73$  & $2\cdot11\cdot13\cdot19\cdot347$        & $21$ & $19$ & $21$ & $+2$\\
$1899206$ & $2\cdot19\cdot23\cdot41\cdot53$          & $2\cdot3\cdot5\cdot29\cdot37\cdot59$    & $21$ & $19$ & $21$ & $+2$\\
$1952786$ & $2\cdot13\cdot19\cdot59\cdot67$          & $2\cdot3\cdot5\cdot7\cdot17\cdot547$    & $21$ & $19$ & $20$ & $+1$\\
$1984814$ & $2\cdot5\cdot41\cdot47\cdot103$          & $2\cdot3\cdot11\cdot17\cdot29\cdot61$   & $21$ & $19$ & $19$ & $+0$\\
\bottomrule
\end{tabular}
\caption{Prime-rich even values of $k$ in $2.5\times10^{5}<k<2\times10^{6}$ chosen to
maximise $W(k):=\omega_{\mathrm{odd}}(k)+2\,\omega_{\mathrm{odd}}(k^2-16)$; for even $k$
the maximum is $W=21$. All are even with $16\nmid k$, so $E_k$ is additive at $2$,
$c(k)=2$ and $\hat v(k)=W(k)-2$. The table lists every even $k$ in the window attaining
$W=21$, i.e.\ the complete set of even maximisers. In every case $v_2(n_k)\ge\hat v(k)$,
with equality at $k=1540136$, $k=1812794$, and $k=1984814$.}
\label{tab:largeW-even}
\end{table}


\clearpage

\bibliographystyle{alpha}
\bibliography{Bibliography}


\end{document}
